%% 
%% Copyright 2007-2025 Elsevier Ltd
%% 
%% This file is part of the 'Elsarticle Bundle'.
%% ---------------------------------------------
%% 

\documentclass[preprint,12pt]{elsarticle}

%% Use the option review to obtain double line spacing
%% \documentclass[authoryear,preprint,review,12pt]{elsarticle}

%% Use the options 1p,twocolumn; 3p; 3p,twocolumn; 5p; or 5p,twocolumn
%% for a journal layout:
%%\documentclass[final,1p,times]{elsarticle}
%% \documentclass[final,1p,times,twocolumn]{elsarticle}
%%\documentclass[final,3p,times]{elsarticle}
%% \documentclass[final,3p,times,twocolumn]{elsarticle}
%% \documentclass[final,5p,times]{elsarticle}
%% \documentclass[final,5p,times,twocolumn]{elsarticle}

%% For including figures, graphicx.sty has been loaded in
%% elsarticle.cls. If you prefer to use the old commands
%% please give \usepackage{epsfig}

\usepackage{amssymb}
\usepackage{amsmath}
\usepackage{amsmath,amssymb,amsfonts}
\usepackage{graphicx,color}
\usepackage{textcomp}
\usepackage{hyperref}
\usepackage{enumerate} 
\usepackage{hhline}
\usepackage[table,xcdraw]{xcolor}
\usepackage{soul} % text
\usepackage{balance} % text
\usepackage[nocomma]{optidef} % math
\usepackage{float} % figure
\usepackage{array} % table
\usepackage{algpseudocode} % algorithm
\usepackage{algorithmicx} % algorithm
\usepackage[ruled, linesnumbered]{algorithm2e} % algorithm
\usepackage{subfig}
\usepackage{makecell}
\usepackage{multirow}
\usepackage{colortbl}
\usepackage[most]{tcolorbox}
\usepackage{url}


\journal{Clinical eHealth}

\begin{document}

\begin{frontmatter}

%% Title, authors and addresses

%% use the tnoteref command within \title for footnotes;
%% use the tnotetext command for theassociated footnote;
%% use the fnref command within \author or \affiliation for footnotes;
%% use the fntext command for theassociated footnote;
%% use the corref command within \author for corresponding author footnotes;
%% use the cortext command for theassociated footnote;
%% use the ead command for the email address,
%% and the form \ead[url] for the home page:
%% \title{Title\tnoteref{label1}}
%% \tnotetext[label1]{}
%% \author{Name\corref{cor1}\fnref{label2}}
%% \ead{email address}
%% \ead[url]{home page}
%% \fntext[label2]{}
%% \cortext[cor1]{}
%% \affiliation{organization={},
%%             addressline={},
%%             city={},
%%             postcode={},
%%             state={},
%%             country={}}
%% \fntext[label3]{}

\title{Agentic AI System for Generating Evidence-Based Second Medical Opinions}

%% use optional labels to link authors explicitly to addresses:
%% \author[label1,label2]{}
%% \affiliation[label1]{organization={},
%%             addressline={},
%%             city={},
%%             postcode={},
%%             state={},
%%             country={}}
%%
%% \affiliation[label2]{organization={},
%%             addressline={},
%%             city={},
%%             postcode={},
%%             state={},
%%             country={}}



\cortext[cor1]{Corresponding author}
\author[1]{Diana Hawashin\corref{cor1}}
\ead{diana.jhawashin@ku.ac.ae}
\author[1]{Khaled Salah}
\ead{khaled.salah@ku.ac.ae}
\author[2]{Raja Jayaraman}
\ead{rjayaram@nmsu.edu}
\author[3]{Samer Ellahham}
\ead{ellahas@clevelandclinicabudhabi.ae}



\address[1]{Department of Computer and Information Engineering, Khalifa University, Abu Dhabi, UAE}
\address[2]{Department of Industrial Engineering, New Mexico State University, NM, USA}
\address[3]{Heart \& Vascular Institute, Cleveland Clinic Abu Dhabi, Abu Dhabi, UAE}


%% Abstract

\begin{abstract}
Second medical opinions play a critical role in minimizing diagnostic uncertainty and improving patient outcomes. However, healthcare systems frequently struggle to provide them in a timely manner due to complex clinical workflows and limited access to up-to-date, evidence-based guidelines. Current AI-based clinical decision support systems are largely restricted to narrow medical domains and rely on static knowledge repositories, limiting their ability to produce reliable, guideline-consistent reasoning across diverse conditions. In this paper, we introduce an agentic AI framework that delivers evidence-based second medical opinions by integrating Retrieval-Augmented Generation (RAG) with a controller-driven orchestration architecture. The framework dynamically retrieves pertinent clinical guidelines through semantic chunking and vector similarity search, and it coordinates multiple large language models (LLMs), including GPT-4o, Claude Opus 4.6, and Gemini 2.5 Flash, to produce structured recommendations aligned with clinical guidelines. A web-based interface allows users to submit patient information and obtain standardized outputs encompassing condition status, diagnosis, severity assessment, and treatment planning. Evaluation across representative respiratory scenarios shows that the agentic pipeline achieves strong reasoning-level guideline agreement under a conservative G-Eval threshold, matching the strongest single-model configuration (precision = 1.00, recall = 0.75, specificity = 1.00, F1-score = 0.86), while maintaining low per-case token-based costs (average USD~0.0473 across four cases), mean end-to-end latency of 44.66 seconds in the tested environment, strong semantic alignment with guideline-derived reference answers (mean cosine alignment 0.7629), and stable repeatability across 30 controlled executions (CV = 2.8\%). These findings indicate that agentic coordination can improve consistency and clinical reliability beyond what is achievable with individual LLM configurations alone.
\end{abstract}





%% Keywords
\begin{keyword}
Agentic Architecture, LLMs, Medical Diagnosis, RAG, Clinical Decision Support Systems
\end{keyword}

\end{frontmatter}

\section{Introduction}
\label{sec:intro}

Second medical opinions play a critical role in reducing diagnostic uncertainty and improving patient safety, particularly in complex or high-risk clinical decisions. As healthcare delivery becomes increasingly specialized and data-intensive, clinicians must synthesize expanding bodies of clinical guidance while maintaining high diagnostic accuracy \cite{r0}. Diagnostic errors remain a significant contributor to preventable harm in healthcare systems, accounting for approximately 46\% of preventable adverse events in primary and emergency care settings \cite{r24}. Prior studies further indicate that obtaining a second opinion can meaningfully alter diagnoses or treatment plans across diverse clinical contexts \cite{r25,r26}.

Despite their recognized value, second-opinion processes often remain inefficient, particularly in high-risk specialties such as oncology and surgical care where diagnostic and therapeutic decisions carry substantial consequences. Clinicians frequently consult lengthy guideline documents, reconcile heterogeneous recommendations, and apply severity thresholds under time pressure. Guideline sources may be distributed across organizations and updated independently, increasing the risk of inconsistent interpretation. Prior research across multiple specialties indicates that structured second-opinion review can improve diagnostic accuracy, strengthen patient understanding of treatment options, and increase confidence in care decisions \cite{new1,new2,new3,new5,new6}. At the system level, electronic consultation platforms have been introduced to facilitate specialist input and improve coordination \cite{new4,new7}. Nevertheless, variability in interpretation and fragmented information workflows persist, which highlights the need for scalable, structured mechanisms to support consistent second-opinion generation.

The promise of large language models (LLMs) has generated interest in automating aspects of clinical reasoning. Systems such as Claude, Gemini, and GPT have demonstrated strong performance on standardized medical examinations and benchmark reasoning tasks \cite{r15,r16}. However, performance on structured tests does not guarantee reliability in real-world clinical environments. Empirical investigations report that LLMs may fail to request appropriate diagnostic tests or fully adhere to established clinical guidelines in complex care settings \cite{r27,r32}. Other studies document risks related to demographic bias and unsupported or hallucinated outputs \cite{r28,r30,r31}. For example, one analysis found that an LLM produced incorrect diagnoses in 72\% of pediatric cases and failed to provide sufficiently specific diagnoses in additional instances when evaluated without structured clinical grounding \cite{r29}. Additional analyses emphasize persistent challenges in domain adaptation and structured clinical reasoning \cite{r33,rodrigues2025}. These findings suggest that unconstrained text generation alone is insufficient for safety-critical clinical support. Improving reliability in medical AI has therefore increasingly focused on grounding model outputs in explicit evidence. Approaches that incorporate structured evidence selection and transparent reasoning mechanisms have demonstrated improved interpretability and factual consistency in healthcare applications \cite{emin2025}. Such strategies emphasize that clinical decision support requires not only accurate language generation, but also controlled access to authoritative sources and structured reasoning pathways.

Retrieval-Augmented Generation (RAG) represents one practical method for grounding LLM outputs in external knowledge. By retrieving relevant documents at inference time, RAG reduces hallucination and allows systems to incorporate updated clinical material without retraining. Within healthcare, RAG has been applied to medical chatbots \cite{r16}, clinical trial eligibility screening \cite{r18}, extraction of risk factors from electronic health records \cite{r17}, and radiology workflows \cite{r39}. However, most implementations remain limited to single-turn question answering or narrowly scoped tasks. They provide limited support for multi-stage reasoning and do not consistently enforce structured guideline alignment throughout a full decision-making workflow. Clinical second-opinion generation, in contrast, is inherently multi-step. It requires identifying relevant evidence, applying condition-specific severity criteria, evaluating treatment pathways, comparing alternatives, and synthesizing recommendations in a structured format. These characteristics align more naturally with orchestrated, tool-based reasoning frameworks than with single-pass language generation.

Agentic AI has emerged as a design paradigm for coordinating such multi-step reasoning processes. Agentic architectures integrate planning, tool invocation, and structured intermediate reasoning to manage complex tasks \cite{agenticintro1,agenticintro2,agenticintro3}. Frameworks such as ReAct \cite{r35} and Least-to-Most prompting \cite{r36} decompose reasoning into explicit stages. Clinical-oriented systems including MedGen \cite{r22}, ClinicalAgent \cite{r20}, and AI Hospital \cite{r21} extend these ideas to healthcare scenarios. While these approaches demonstrate advances in reasoning structure and coordination, many operate within restricted domains or rely on fixed knowledge sources without dynamic guideline retrieval.

In this work, we design and evaluate an agentic AI system for generating evidence-based second medical opinions. The system integrates dynamic guideline retrieval with controller-based orchestration of multiple LLM tools to produce structured, guideline-aligned recommendations. Evaluation is conducted within respiratory medicine using authoritative WHO and MSF guideline sources as a controlled domain, while the architectural design itself is not restricted to this specialty. The system functions strictly as a clinical decision-support tool and does not replace clinician judgment or perform autonomous diagnosis. Our key contributions are as follows:

\begin{itemize}

\item We propose an agentic AI framework for generating structured second medical opinions that supports both new and follow-up cases within the evaluated respiratory domain, combining real-time guideline retrieval with coordinated multi-LLM reasoning.

\item We introduce a controller-based architecture that manages retrieval, dynamic tool invocation, and structured synthesis to enhance reasoning consistency and guideline alignment.

\item We evaluate the framework against baseline and single-model RAG configurations, reporting strong reasoning-level guideline agreement under a conservative G-Eval threshold (precision \(=1.00\), recall \(=0.75\), specificity \(=1.00\), and F1-score \(=0.86\) for the agentic configuration), strong semantic alignment with guideline-derived reference answers (mean cosine alignment 0.7629), low per-case token-based costs (average USD~0.0473 across four cases with dynamic tool usage), and mean end-to-end response times of 44.66 seconds within the tested environment.

\item We assess output stability across repeated runs to examine consistency and operational feasibility within the evaluated scenarios, yielding stable repeatability across 30 executions (CV = 2.8\%).

\end{itemize}

The remainder of this paper is structured as follows. Section \ref{sec:related_work} reviews prior work on clinical decision support systems (CDSS), RAG, and agentic AI. Section \ref{sec:methodology} presents the research methodology adopted in this paper. Section \ref{sec:system_architecture} describes the overall system architecture, including controller design and component interactions. Section \ref{sec:implementation} details the implementation of individual modules and tool selection. Section \ref{sec:test} presents evaluation results, including guideline adherence, classification performance, latency, and output consistency. Section \ref{sec:discussion} analyzes key findings and their implications for clinical deployment. Section \ref{sec:conclusion} summarizes the contributions and outlines directions for future work.

\section{Related Work}
\label{sec:related_work}

In this section, we review prior work on medical decision support systems, RAG in clinical settings, and multi-agent reasoning frameworks. The discussion highlights limitations in existing approaches related to guideline integration, domain coverage, and system reliability, and identifies the specific gaps that motivate the design of our proposed system.


\subsection{Medical Decision Support Systems}

Clinical decision support systems are designed to improve diagnostic accuracy, reduce variability in care, and support clinicians in complex decision-making.

The authors in \cite{r9} review the evolution of CDSS and highlight their integration with electronic medical records to support complex decisions. They identify issues such as alert fatigue, workflow disruptions, and unintended consequences, and recommend evidence-based design and thorough evaluation as necessary improvements. Similarly, the authors in \cite{r13} survey recent AI-based CDSS and confirm their potential to improve diagnostic accuracy and reduce errors, however, they highlight inconsistent evidence quality and emphasize the need for robust real-world validation.

In \cite{r10}, the authors propose an AI-driven mental health assessment system using a 28-question tool powered by the NEPAR algorithm and machine learning models trained on participant and historical data. The system achieves 89\% diagnostic accuracy and aims to reduce clinician workload through automated diagnosis.
The ERVIN.2 system presented by \cite{r11} supports joint decision-making for osteoarthritis patients considering hip or knee replacement. It displays predicted one-year outcomes for surgical and non-surgical treatments using real-time patient-reported data.

In \cite{r12}, the authors introduce a deep learning-based CDSS for heart disease prediction using dense neural network architectures. Their model achieves high accuracy, sensitivity, and specificity across four benchmark datasets.

The study by \cite{r14} examines the effects of explainable AI (XAI) on trust calibration in CDSS through a prescription-screening tool tested with 41 medical practitioners. They find that while explanations improved overall Human-AI performance, they also increase over-reliance on the system and fail to help clinicians detect incorrect recommendations.

Taken as a category, traditional and AI-based CDSS have demonstrated effectiveness in task-specific or disease-specific applications. However, most operate within predefined decision rules or trained predictive models that do not dynamically retrieve evolving guideline documents at inference time. Consequently, their adaptability to multi-condition reasoning scenarios remains limited when compared to systems that explicitly integrate real-time evidence retrieval.

Unlike conventional CDSS, which primarily rely on predefined rule sets or trained predictive models embedded within electronic record systems, our proposed framework dynamically retrieves external guideline documents at inference time and coordinates multiple reasoning agents to synthesize structured outputs. This distinction shifts the system from static decision assistance toward adaptive, evidence-grounded orchestration across diverse clinical scenarios.



\subsection{RAG for Clinical NLP}

Retrieval-Augmented Generation (RAG) has been applied in clinical natural language processing (NLP) to improve factual grounding and evidence linkage in generative systems. By incorporating external clinical sources at inference time, RAG-based models can provide context-aware and document-grounded outputs. Within healthcare, RAG approaches can be broadly categorized into classical relevance-based retrieval systems, clinically tuned RAG pipelines, and emerging guideline-oriented retrieval frameworks.


Classical RAG implementations focus primarily on semantic similarity matching between user queries and document collections. In clinical settings, such systems have been used to enhance chatbot reliability and document-grounded responses. The authors in \cite{r15} explore the use of RAG to improve chatbot responses in clinical decision-making by integrating current clinical data and reliable guidelines. Similarly, the authors in \cite{r16} develop MEDGPT, a healthcare chatbot that retrieves content from PDFs and PubMed to generate context-relevant responses. These systems emphasize contextual relevance and factual grounding, but typically rely on static retrievers optimized for similarity rather than structured enforcement of clinical protocols.

Clinically tuned RAG systems extend this paradigm by applying retrieval mechanisms to domain-specific tasks and structured clinical corpora. In \cite{r17}, RAG-enabled generative AI is used to summarize nutrition-related risk factors from EHR notes in aged care facilities in Australia, achieving 99\% summarization accuracy and 90\% extraction accuracy. The RECTIFIER system introduced in \cite{r18} applies GPT-4 with RAG to extract clinical trial eligibility criteria from unstructured notes, demonstrating high sensitivity and specificity compared to expert reviews. Moreover, the authors in \cite{r19} use a RAG-enabled LLM to identify major bleeding events from EHR data for research and quality assessment, achieving 72\% precision and recall by combining document retrieval with chain-of-thought prompting. These systems demonstrate improved domain adaptation and factual precision, yet they remain largely focused on single-task or single-turn inference rather than coordinated multi-stage clinical reasoning.

More recent discussions emphasize the importance of aligning retrieval mechanisms with authoritative medical guidance. However, most existing RAG-based clinical systems do not explicitly enforce structured guideline adherence across sequential reasoning stages within a comprehensive decision-support workflow.

From a structural perspective, RAG-based approaches improve factual grounding by linking generation to retrieved documents, particularly in extraction and summarization tasks. Yet most implementations emphasize isolated question answering or document-level inference rather than coordinated multi-step reasoning across clinical dimensions. Their retrieval mechanisms are typically optimized for semantic relevance rather than structured enforcement of clinical guideline constraints within a complete decision-support workflow.

While RAG systems enhance factual grounding through document retrieval, they typically operate within a single-model generation pipeline. In contrast, our approach integrates retrieval with explicit orchestration across multiple LLM tools, which enables structured synthesis of intermediate reasoning steps rather than isolated answer generation. The emphasis is therefore not only on grounding, but on coordinated multi-stage clinical reasoning aligned with formal guideline structures.



\subsection{Multi-Agent and Ensemble Reasoning in Clinical AI}

Multi-agent reasoning frameworks aim to address the limitations of single-LLM systems by decomposing tasks, assigning roles, and coordinating multiple agents to perform specialized functions. These architectures show promise in domains such as diagnostic conversations and clinical trial prediction.

ClinicalAgent \cite{r20} applies multi-agent reasoning to clinical trial outcome prediction by combining GPT-4 with Least-to-Most and ReAct strategies, achieving a 0.33 improvement over standard prompting with a PR-AUC of 0.7908. The system demonstrates the value of structured reasoning decomposition within a predefined prediction task. However, its architecture is optimized for a specific trial-classification objective rather than dynamically retrieving and enforcing evolving clinical guideline documents during inference.

AI Hospital \cite{r21} simulates multi-turn diagnostic dialogue through GPT-4-based agents representing roles such as Chief Physician, Examiner, and Patient. Its dispute resolution mechanism enhances internal consistency across simulated agents. While this approach supports conversational diagnostic reasoning, it primarily operates within a simulated role-based framework rather than integrating structured retrieval of external clinical protocols as a governing constraint.

MedGen \cite{r22} structures clinical reasoning into goal setting, data collection, argument linking, and plan selection, with demonstrations in depression and breast cancer scenarios. The framework emphasizes interpretability through structured outputs and reasoning transparency. Nevertheless, it relies on predefined reasoning structures and domain-bounded knowledge rather than embedding dynamic retrieval of authoritative guideline repositories within each reasoning stage.

MDAgents and Health-LLM \cite{r23} extend multi-agent coordination to collaborative decision-making and multimodal prediction, integrating wearable sensor data with contextual metadata. These systems illustrate how coordinated agents can enhance predictive performance across heterogeneous inputs. However, their focus remains on multimodal modeling and collaborative inference rather than structured alignment with externally retrieved clinical guideline documents.

From an architectural perspective, multi-agent and ensemble systems improve reasoning diversity and intermediate structure by distributing subtasks across coordinated agents. Yet in most existing implementations, coordination mechanisms emphasize reasoning decomposition within internally available knowledge or static corpora. Explicit integration of real-time guideline retrieval as a core orchestration component remains limited.

Among the approaches most closely related to our work are ClinicalAgent \cite{r20}, MedGen \cite{r22}, and AI Hospital \cite{r21}, which incorporate multi-agent reasoning and structured intermediate steps in healthcare contexts. In contrast, our proposed framework embeds dynamic guideline retrieval directly within the controller layer and synthesizes outputs across multiple LLM tools into structured second medical opinions grounded in authoritative protocol documents. This integration of real-time evidence retrieval with coordinated multi-model orchestration distinguishes the present system from prior multi-agent clinical AI architectures that prioritize reasoning decomposition without structured guideline enforcement.

\subsection{Research Gap}

Despite notable advances in clinical AI, existing decision support systems have not achieved robust, general-purpose reasoning capabilities suitable for real-world healthcare environments. Many systems remain confined to narrow clinical domains or single-condition tasks and do not consistently produce outputs aligned with the structured reasoning processes expected in expert medical decision-making. As a result, their recommendations may lack generalizability, transparency, and protocol-level grounding.

RAG-based methods improve evidence linkage by incorporating external sources at inference time, yet most implementations remain optimized for single-turn question answering or predefined extraction tasks. They do not typically coordinate multi-stage reasoning processes or enforce structured alignment with evolving clinical guideline documents across an end-to-end decision workflow. Similarly, multi-agent and ensemble LLM architectures enhance reasoning decomposition and role-based coordination, but often operate over static corpora or domain-bounded knowledge representations without embedding real-time guideline retrieval as a governing constraint within the orchestration layer.

Consequently, an architectural gap persists where prior systems have not unified dynamic retrieval of authoritative clinical guideline repositories with controller-based multi-LLM orchestration and structured synthesis of second medical opinions within a single framework. In contrast, our approach integrates real-time guideline retrieval directly into the orchestration process and coordinates multiple LLM tools to generate structured, guideline-aligned recommendations across multi-condition scenarios. This enables protocol-grounded, evidence-aware reasoning while maintaining structured output consistency appropriate for clinical decision support.

\section[]{Research Methodology}
\label{sec:methodology}

This study adopts the Design Science Research Methodology (DSRM) proposed by Peffers et al. \cite{Peffers01122007} to guide the development and assessment of the proposed agentic AI system. The research centers on the construction of a controller-based architecture that integrates dynamic guideline retrieval with coordinated multi-LLM orchestration, and on the systematic evaluation of this artifact against defined functional and operational objectives. The DSRM phases are operationalized in relation to the specific components of the system, the respiratory-domain instantiation used for controlled demonstration, and the quantitative metrics applied for performance evaluation. The methodological progression follows problem identification, objective definition, artifact design and development, domain-specific demonstration, quantitative evaluation, and knowledge dissemination, as illustrated in Figure \ref{fig:dsr_process}.


\begin{figure*}[t]
\centering
\includegraphics[width=\linewidth]{DSRM.png}
\caption{Design Science Research Methodology (DSRM) applied to the proposed agentic AI system for second medical opinion generation, illustrating the relationship between research objectives, system design, respiratory-domain instantiation, and quantitative evaluation.}
\label{fig:dsr_process}
\end{figure*}

\subsection[]{Problem Identification and Motivation}

The research is motivated by limited access to timely, evidence-based second medical opinions in routine clinical practice. Clinicians often must synthesize guidelines under time pressure, which risks inconsistent application. The study frames this gap as a design challenge: develop an artifact that dynamically retrieves authoritative clinical guidance and produces structured, guideline-aligned second opinions to support clinician decision-making.

\subsection[]{Objectives of a Solution}

From the identified problem we derive four evaluation objectives for the artifact: (1) dynamically retrieve relevant guideline material at inference time; (2) coordinate multiple LLMs to produce consistent reasoning; (3) synthesize structured, traceable second opinions for new and follow-up cases; and (4) meet operational constraints (cost, latency, output stability). These objectives define the success criteria for subsequent demonstration and evaluation.



\subsection[]{Design and Development}

The artifact is an agentic AI system that integrates a guideline knowledge base and retriever with an orchestration agent that coordinates multiple LLMs and a response synthesizer. Design choices prioritized modularity, reproducibility, and alignment with clinical reasoning workflows rather than autonomous diagnosis. The complete system architecture and component interactions are described in Section \ref{sec:system_architecture}.

\subsection[]{Demonstration}

In the DSRM framework, demonstration establishes that the artifact can be instantiated and executed within a relevant problem context. In this study, the agentic system is instantiated within respiratory medicine as a bounded clinical domain. This domain was selected because it offers publicly accessible and well-structured guideline repositories, including WHO and MSF protocols, that define explicit diagnostic criteria, severity classifications, and treatment pathways. Respiratory conditions such as asthma and pneumonia also involve multi-factor reasoning across age groups and follow-up scenarios, making them suitable for evaluating structured second-opinion generation.

Representative simulated cases, including both new and follow-up consultations, are used to verify functional integration of guideline retrieval, multi-LLM orchestration, and structured opinion synthesis. This stage confirms that the complete reasoning pipeline operates coherently under controlled conditions.



\subsection[]{Evaluation}

Evaluation assesses the extent to which the instantiated artifact satisfies the defined research objectives. The assessment is structured around the four objectives introduced in Section \ref{sec:methodology}. Objective 1, dynamic retrieval of authoritative clinical guidelines, is evaluated through guideline alignment scoring using cosine similarity and G-Eval comparisons against reference answers. Objective 2, coordinated multi-LLM reasoning, is examined through comparative performance analysis across baseline, single-model RAG, and agentic configurations using precision, recall, specificity, and F1 metrics. Objective 3, structured and traceable second opinion generation, is assessed through inspection of standardized output components and consistency across new and follow-up scenarios. Objective 4, operational feasibility, is evaluated through cost analysis, latency measurement, and repeated-run stability testing.

Results and comparative analyses are presented in Sections \ref{sec:test} and \ref{sec:discussion}, where each objective is examined quantitatively and qualitatively within the respiratory-domain evaluation setting.



\subsection[]{Methodological Limitations}

This study follows early-stage design science research and therefore carries methodological constraints consistent with controlled artifact evaluation. The system is assessed using simulated clinical scenarios within a respiratory-domain setting rather than through deployment in live clinical environments. Although the test cases are constructed to reflect realistic diagnostic contexts, they do not represent full-spectrum patient complexity or real-world variability captured in electronic health records.

The evaluation does not include direct clinician-in-the-loop assessment or independent expert validation of generated second opinions. Guideline alignment is measured using structured reference comparisons and automated scoring metrics, which may not fully capture inter-clinician interpretive differences.

Furthermore, the current implementation operates independently of integrated EHR systems. Patient inputs are structured manually within the experimental interface, and the system does not access longitudinal clinical records or institutional data infrastructures. These limitations indicate that further validation in real-world clinical settings, including clinician-based evaluation and EHR integration, is necessary to establish external validity and deployment readiness.




\begin{figure*}[t]
\centering
\includegraphics[width=\textwidth]{SystemArchitecture.png}
\caption{System architecture of the Agentic AI system for second medical opinion generation}
\label{fig:system_architecture}

\end{figure*} 

\section{System Design}
\label{sec:system_architecture}

In this section, we describe the overall architecture of the agentic AI system for generating evidence-based second medical opinions and outline its main components and their interactions.

\subsection{System Overview}

Figure \ref{fig:system_architecture} provides an overview of the proposed agentic AI system, which links a guideline knowledge base and retrieval pipeline with an agentic reasoning and decision-support layer. The architecture comprises a Guideline Knowledge Base and Retrieval module and an Agentic Reasoning and Decision Support module. Together, these modules transform structured clinical guideline documents and patient inputs into unified, guideline-aligned second opinions.

The Guideline Knowledge Base and Retrieval module manages the ingestion of official clinical guidelines provided as text-based PDF documents. Documents are divided into semantically coherent chunks, embedded, and stored in a vector database together with source metadata. This structure enables the retrieval layer to return guideline excerpts specific to the clinical context of each case. Follow-up continuity is handled separately through local visit-state lookup rather than through the vector index itself.

The Agentic Reasoning and Decision Support module operates on top of this retrieval layer. When a clinician submits a structured patient case through the user interface, the orchestration agent determines the visit type, retrieves relevant guideline chunks, checks for prior local visit context when the case is a follow-up, and constructs structured prompts for one or more expert LLMs. When multiple tools are invoked, each consulted LLM produces an independent clinical recommendation conditioned on retrieved guideline context. The controller then synthesizes the available tool outputs into a single structured report that is returned to the clinician. The detailed interaction flow among user, orchestrator, retriever, LLM tools, and synthesizer is illustrated in Figure \ref{fig:seq}.

A controller-based agentic design was selected to introduce explicit coordination across retrieval, reasoning, and synthesis stages. In conventional single-LLM RAG configurations, retrieved guideline excerpts are injected into a single generation step, which leaves reasoning structure and output consistency dependent on one model instance. In contrast, the controller manages retrieval and routing decisions, invokes one or more expert tools as needed, and enforces a structured response schema during final synthesis.

This orchestration layer enables comparison across independently generated expert responses and reduces reliance on a single generative trajectory. When divergence occurs or retrieved evidence is sparse, the controller synthesizes across the available tool outputs conditioned on the retrieved guideline excerpts to produce a conservative unified recommendation. By separating retrieval management, expert consultation, and response synthesis into coordinated stages, the architecture supports more consistent guideline alignment than single-model RAG pipelines that lack explicit reasoning control.




\subsection{System Components}

The system consists of modular components that span knowledge ingestion, semantic retrieval, agentic reasoning, and interface delivery. Together, these components form a pipeline that processes guideline documents and patient data into structured, evidence-based recommendations. Below, we describe each component in the order it appears in the system workflow.

\begin{itemize}

\item \textbf{Data Source Files}: This component represents official medical guideline documents provided to the system as machine-readable PDF files. These documents serve as the knowledge foundation for the system and ensure that clinical recommendations follow trusted evidence-based standards. 
In the current implementation, guideline ingestion occurs through local provision of text-based PDF documents, which are subsequently parsed, chunked, and embedded into the vector database. Updates to guideline content require reprocessing and re-indexing of the corresponding documents; incremental synchronization with external repositories is not performed automatically. When multiple guideline sources such as WHO and MSF are included, they are indexed within the same retrieval space and differentiated through metadata tags. The system does not algorithmically resolve conflicts between sources; instead, semantic retrieval returns relevant excerpts from available documents, and any differences are reflected in the retrieved context provided to the reasoning layer.


\item \textbf{Chunking Process}: The chunking module divides long guideline documents into smaller, overlapping segments. Each chunk preserves semantic coherence while remaining within model input constraints. Overlapping windows maintain context at segment boundaries. As shown in \cite{r5}, this method supports efficient document retrieval by balancing granularity and contextual integrity.

\item \textbf{Embedding Model}: After chunking, the system uses an embedding model to convert text segments into dense vector representations. These embeddings support semantic similarity search and enable the retrieval of guideline content that aligns with patient cases, even when phrasing varies. Dense retrieval has been shown to improve performance in medical question answering and other NLP tasks \cite{r6}.

\item \textbf{Vector Store}: This component stores the vector embeddings of guideline chunks together with lightweight source metadata. It supports high-dimensional similarity search using nearest-neighbor techniques, allowing the system to retrieve guideline content that is relevant to clinical queries. Vector databases play a critical role in enabling scalable, semantically meaningful retrieval in generative systems \cite{r7}.

\item \textbf{Visit State Handler}: Follow-up continuity is managed through a local visit-state mechanism keyed by patient identifier. This component retrieves the most recent prior visit record when a follow-up case is submitted, allowing the controller to compare current and prior clinical context without storing those records in the FAISS retrieval index.

\item \textbf{Retriever}: The retriever queries the vector store to identify the most relevant guideline chunks for a given patient case. It filters and ranks these results before passing them to the orchestration agent, ensuring that retrieved content aligns with the case's clinical context.

\item \textbf{Orchestration Agent}: This controller module coordinates the reasoning pipeline. It receives structured patient case data, determines the type of visit, manages the retrieval process, selects appropriate LLM tools, and integrates their outputs into a unified recommendation.

\item \textbf{Large Language Models}: These LLMs receive structured prompts containing patient case details and guideline excerpts. Each model returns an independent recommendation that reflects expert-level reasoning based on the retrieved evidence.

\item \textbf{Response Synthesizer}: The synthesizer combines multiple LLM responses into a single, consistent clinical recommendation. It uses structured merging techniques and follows self-consistency principles to reduce output variance and enhance reliability \cite{r8}.

\item \textbf{User Interface}: The interface allows clinicians to submit new or follow-up cases and review structured outputs that include condition classification, diagnosis, severity, and treatment planning. It ensures consistent data collection and accessible decision support.

\end{itemize}

The system assumes that guideline documents are provided in machine-readable text form and that patient data are entered as structured textual fields. It does not process image-based guideline content, handwritten notes, or raw electronic health record databases, and it operates independently of institutional clinical information systems.



\subsection{System Workflow and Agentic Reasoning Sequence}

\begin{figure*}[t]
\centering
\includegraphics[width=\textwidth]{seq22.png}
\caption{Sequence diagram of the Agentic AI pipeline for medical recommendation generation}
\label{fig:seq}
\end{figure*} 

Figure~\ref{fig:seq} illustrates the step-by-step interaction flow of the agentic AI system during second medical opinion generation. The system processes structured patient data and coordinates retrieval, reasoning, and synthesis to produce a guideline-aligned recommendation.

The process begins when a clinician enters patient case information through the interface. The interface standardizes this input and forwards it to the Agentic Orchestrator, which identifies whether the visit is new or a follow-up. For new visits, the orchestrator queries the vector store to retrieve guideline passages that match the submitted case details. For follow-up visits, it also retrieves the most recent prior visit record via the local visit-state mechanism to support longitudinal comparison and updated context alignment.

Next, the orchestrator assesses case complexity to select the appropriate expert LLMs. Based on the perceived complexity and risk of the case, it may query one, two, or all available models. Each model receives a structured prompt that includes patient details and retrieved guideline context, and returns a structured recommendation.

The orchestrator then integrates the outputs using a synthesis strategy that emphasizes consistency and guideline adherence. When multiple expert LLM tools are consulted, their responses serve as intermediate reasoning steps within the agentic framework. The controller incorporates these tool outputs into its reasoning trajectory and produces a single structured recommendation based on the combined guideline-grounded evidence. Rather than applying explicit majority voting or numeric thresholds, the synthesis relies on the controller's structured reasoning over retrieved guideline context and tool responses to generate a unified opinion. The unified recommendation is returned through the interface as a structured report that includes diagnosis, severity, and treatment plan. This output format ensures transparency, reduces variability, and supports reliable decision-making aligned with established clinical protocols.




\begin{table}[!ht]
\centering
\caption{System Setup and Environment Summary}
\label{tab:setup_summary}
\renewcommand{\arraystretch}{1.2}
\begin{tabular}{|l|p{8cm}|}
\hline
\textbf{Item} & \textbf{Configuration} \\
\hline
Execution Environment & Python-based application with hosted LLM APIs and local FAISS retrieval \\
\hline
Programming Language & Python 3 \\
\hline
Embedding Model & OpenAI \texttt{text-embedding-ada-002} (1536-dimensional) \\
\hline
Vector Store & FAISS \\
\hline
Chunking Method & Overlapping window, 800 characters per chunk with 100-character overlap \\
\hline
Text Extraction & PyPDFLoader \\
\hline
Chunk Splitter & RecursiveCharacterTextSplitter \\
\hline
LLM APIs Used & GPT-4o (OpenAI), Claude Opus 4.6 (Anthropic), Gemini 2.5 Flash (Google) \\
\hline
Temperature Setting & 0.2 (applied uniformly) \\
\hline
Orchestration Framework & Python controller logic with LangChain model wrappers \\
\hline
Environment Variable Handling & \texttt{dotenv} with runtime key checks \\
\hline
\end{tabular}
\end{table}

\section{System Implementation}
\label{sec:implementation}
In this section, we describe the practical setup of our system, which integrates an agentic multi-LLM pipeline with retrieval-augmented reasoning to ensure outputs remain dynamically aligned with relevant clinical guidelines. The implementation includes document embedding and chunking for guideline ingestion, semantic retrieval components, custom prompt engineering for each LLM, a controller-based orchestration step for structured synthesis, and an interactive web interface for clinical use. The released source code, application implementation, and sample study inputs are available in our public GitHub repository\footnote{\url{https://github.com/djh97/smo}}.



\subsection{Setup and Environment}
\label{sec:setup_environment}

The system is implemented in Python using a modular architecture that supports reproducible experimentation and controlled evaluation. The runtime environment enables integration of retrieval components, multi-LLM orchestration, and interface delivery within a unified workflow. API keys for external LLM services are loaded securely at runtime, and the released research prototype supports optional local persistence of structured visit outputs to enable follow-up review and reproducibility. Table~\ref{tab:setup_summary} summarizes the primary configuration parameters relevant to retrieval, orchestration, and model execution.

All patient case data used in this study consist of structured, simulated clinical scenarios rather than real patient records. During execution, case data are transmitted to external LLM APIs over encrypted HTTPS connections. No personally identifiable health information is included in the experimental cases, and no patient data are stored in external databases. Any locally saved visit outputs in the released prototype are limited to the structured simulated cases used for experimentation.


\subsection{System Components and Configuration Rationale}
\label{sec:component_selection}

This subsection outlines the core components and configuration strategies that underpin the system's functionality. Each element is selected to support transparent, dynamic, and guideline-aligned clinical reasoning within a modular, cloud-compatible architecture.


\subsubsection{Orchestration Agent}
The system uses LangChain model wrappers within a custom Python controller to implement the orchestration layer that manages multiple LLM tools and coordinates agentic reasoning. These wrappers support structured model invocation and consistent context handling across providers, while the custom controller implements routing and synthesis logic for the retrieval-augmented multi-LLM workflow.

\subsubsection{Chunking Process}

Guideline documents are segmented using an overlapping sliding-window strategy with a chunk size of 800 characters and an overlap of 100 characters.  This configuration balances retrieval granularity and contextual continuity while remaining compatible with LLM input constraints. The chunking logic is defined as:

\begin{equation}
C_i = D[s_i : s_i + L]
\end{equation}

\begin{equation}
s_{i+1} = s_i + (L - O)
\end{equation}

where \(D\) denotes the document text, \(L\) the chunk length, and \(O\) the overlap parameter.


\subsubsection{Embedding Model}
The system employs the OpenAI \texttt{text-embedding-ada-002} model to generate 1536-dimensional vector representations for guideline segments and structured patient inputs. These embeddings are stored in the FAISS index to support semantic similarity search during retrieval. This embedding model was selected due to its demonstrated performance in semantic retrieval tasks and compatibility with large-scale vector indexing. Its dimensionality provides sufficient representational capacity for capturing clinically relevant context while maintaining efficient similarity computation.


\subsubsection{Vector Store}
FAISS is used as the vector storage and retrieval backend. For each patient case, the system retrieves the top \(k = 3\) guideline chunks using semantic nearest-neighbor search between the embedded case representation and stored guideline vectors. The selection of \(k = 3\) reflects a balance between contextual completeness and prompt efficiency. Retrieving too few segments risks omitting relevant guideline criteria, whereas retrieving excessive segments increases token load and may dilute focus. This configuration supports efficient nearest-neighbor search over high-dimensional embeddings.



\subsubsection{Text Extractor and Splitter}
The system uses PyPDFLoader to extract text from medical PDF documents, which often contain complex structures and annotations. This tool ensures that the extracted content preserves logical ordering and readability. The RecursiveCharacterTextSplitter then segments the text into semantically coherent overlapping chunks, enabling context-aware RAG retrieval and downstream clinical reasoning.

\subsubsection{LLM Tools}
\begin{itemize}
    \item \textbf{GPT-4o via ChatOpenAI:} Acts as the system's primary reasoning agent due to its strong alignment capabilities and reliable structured output formatting.
    \item \textbf{Claude Opus 4.6 via ChatAnthropic:} Contributes diverse clinical reasoning perspectives and safety-tuned language behavior, increasing ensemble robustness.
    \item \textbf{Gemini 2.5 Flash via ChatGoogleGenerativeAI:} Adds rapid inference capabilities and distinct reasoning styles, improving response diversity and coverage.
\end{itemize}

\subsubsection{Temperature Setting}
The system sets a uniform temperature of 0.2 across all model calls. This relatively low temperature was selected to promote deterministic and stable outputs in safety-critical clinical reasoning tasks, while still allowing limited flexibility for nuanced case differentiation. Higher temperatures were observed during preliminary testing to increase variability without improving guideline alignment.


\subsubsection{Orchestrator Configuration}
The orchestration logic is implemented in Python using LangChain model wrappers together with a custom prompt-driven controller. This design enables dynamic tool selection, structured prompting, and multi-step reasoning. The controller compares outputs across models, synthesizes the result, and adapts its prompting behavior to case complexity. This configuration improves interpretability and robustness while keeping the routing and synthesis logic explicit and auditable.


\subsection[]{Agentic AI Logic and Algorithms}
\label{sec:system_workflow}

This subsection presents the core logic and algorithmic structure of the agentic AI system. It outlines the sequential steps for handling patient input, retrieving relevant clinical guidelines, orchestrating multiple expert LLMs, and synthesizing structured medical recommendations. Each stage is represented through targeted algorithms that illustrate how the system achieves dynamic, guideline-aligned second medical opinions.

\begin{algorithm}[!ht]
\caption{Patient Input, State Initialization, and Guideline-Based Retrieval}
\label{alg:input_and_retrieval}
\SetAlgoLined
\KwIn{Patient data: age, weight, symptoms, $\mathrm{SpO_2}$, heart rate, medical history, visit type, optional patient ID; top-$k$ retrieval parameter}
\KwOut{Structured case dictionary with visit metadata and top-$k$ relevant guideline chunks}
\Begin{
\If{patient ID is not provided or visit type is New}{
Generate new UUID as patient ID
}
\Else{
Use provided patient ID
}

\If{visit type is Follow-up}{
    Retrieve previous visit record from local visit state or saved visit output\\
    Preserve prior data for comparison\
}

Create timestamp for current visit\\

Assemble structured case dictionary using collected inputs\\


Query FAISS using the assembled case text\\

Initialize empty set for duplicates\\
Initialize empty list for filtered chunks\\

\ForEach{retrieved chunk}{
    \If{chunk content not in duplicate set}{
        Append chunk to filtered list\\
        Add chunk content to duplicate set\\
    }
}

Select top-$k$ filtered chunks as guideline context\

\Return case dictionary and guideline chunks\


}
\end{algorithm}



\subsubsection[]{Patient Input Collection, State Handling, and Dynamic Guideline Retrieval}

The system first standardizes structured patient input and initializes visit-level metadata, including patient identifiers and timestamps. For follow-up cases, prior visit context is retrieved to enable longitudinal comparison. The processed case representation is then used to query the vector store for semantically relevant guideline segments using dense similarity search \cite{r6}. The input handling and retrieval procedure is summarized in Algorithm~\ref{alg:input_and_retrieval}.

Continuity across new and follow-up visits is maintained through a patient identifier mechanism. When a new case is created, a UUID-based patient identifier is generated and associated with the structured visit record. Subsequent follow-up visits reference this identifier to retrieve prior case data from the application's local state and, in the released research prototype, from locally saved visit outputs. The study itself uses only structured simulated cases and does not write patient data to external databases. Thus, longitudinal comparison is supported within a controlled research workflow while avoiding use of personally identifiable health information.




\begin{algorithm}[!ht]
\caption{Prompt Construction and Role Encoding}
\label{alg:prompt_construction}
\SetAlgoLined
\KwIn{Case dictionary $C$, guideline chunks $G$, visit type $V$}
\KwOut{Structured prompt $P$}
\Begin{
Initialize prompt template $T$ with role: ``You are a medical agent''\

\If{$V$ is Follow-up}{
    Append prior visit summary to $C$\\
    Add comparison instruction to $T$\
}

Insert $C$ into structured section of $T$\

Insert $G$ into guideline reference section of $T$\

Append output format specification\

\Return formatted prompt $P$\


}
\end{algorithm}


\subsubsection[]{Prompt Construction and Role Encoding}

After retrieval, the system constructs a structured prompt that integrates current case data, optional prior visit context, and selected guideline excerpts. Role instructions and output format constraints are embedded to promote consistent reasoning and reduce generation variability, following established best practices in clinical prompt design \cite{r38,r37}. Synthesis across multiple LLM tools is performed through the agentic reasoning layer rather than through explicit rule-based voting or numerical scoring. Each tool invocation produces a structured response conditioned on the same retrieved guideline context. These responses are incorporated into the controller's reasoning trajectory, and the controller generates a unified structured opinion based on the combined intermediate outputs. The aggregation process therefore reflects coordinated reasoning over multiple expert perspectives, rather than majority voting or trust-weighted averaging. The prompt construction logic is outlined in Algorithm~\ref{alg:prompt_construction}.



\subsubsection[]{Agentic Orchestration, Output Synthesis, and Delivery}

The orchestration stage coordinates multiple expert LLM tools, collects their structured responses, and synthesizes a unified clinical opinion. The controller determines tool invocation and consolidates outputs into predefined clinical components. The orchestration and synthesis process is described in Algorithm~\ref{alg:orchestration_and_output}.


\begin{algorithm}[!ht]
\caption{Agentic Orchestration, Structured Output Generation, and Delivery}
\label{alg:orchestration_and_output}
\SetAlgoLined
\KwIn{Prompt $P$, tool set $\mathcal{T}$}
\KwOut{Final structured clinical opinion}
\Begin{
Initialize empty output list $\mathcal{O}$\;

Select subset of tools $\mathcal{S} \subseteq \mathcal{T}$ based on case needs\;

\ForEach{tool $T_i \in \mathcal{S}$}{
    Submit $P$ to $T_i$\;
    Receive output $O_i$\;
    Append $O_i$ to $\mathcal{O}$\;
}

Construct controller synthesis prompt using patient case, retrieved guideline context, and $\mathcal{O}$\;

Invoke controller LLM once to generate unified 4-point clinical opinion\;

Format output in 4-point structure\;

Display structured output to user interface\;

\Return finalized clinical opinion\;
}
\end{algorithm}


\subsection{Front-End Interface for Clinical Interaction}

The user interface is implemented as a web-based application and serves as the main interaction point for clinical users. It allows structured case submission, triggers the agentic reasoning pipeline, and displays the system's final output in a clinically interpretable format.

\subsubsection{User Input Panel}
The interface collects structured patient data through clearly labeled fields. Inputs include patient ID (optional), visit type (new or follow-up), age, weight (kg), $\mathrm{SpO_2}$ (\%), heart rate (bpm), medical history, and reported symptoms. Input validation ensures all required fields are completed before submission. This structure ensures that the system receives consistent, complete clinical context for generating guideline-aligned second medical opinions.

\subsubsection{Model Interaction and Output}
After the user submits the patient data, the interface triggers the agentic pipeline. This includes dynamic guideline retrieval, prompt construction, and LLM orchestration. Once the reasoning is complete, the system displays the result as structured output panels headed by the final agentic opinion. The response includes four structured components: condition status, likely diagnosis, severity classification, and recommended treatment plan.

\subsubsection{Output Display}
\label{output-display}
The final output appears as structured, read-only presentation panels with consistent formatting. Each of the four clinical components is presented in a dedicated section, enhancing interpretability for practitioners and allowing quick visual parsing for evaluation and follow-up.

\section{Testing and Validation}
\label{sec:test}

This section presents functional testing of the multi-LLM framework using representative clinical cases. The testing covers system execution and compares output quality across standard, RAG-enhanced, and agentic responses.

\subsection{Testing Setup and Clinical Test Cases}
\label{sec:testing_setup}

To validate the system's functionality in a realistic clinical context, we focused on respiratory diseases as the target domain. This selection was informed by the high prevalence of respiratory conditions in global health and the availability of open-access, evidence-based treatment protocols. Two authoritative sources were used: the M\'edecins Sans Fronti\`eres (MSF) Clinical Guidelines and the World Health Organization (WHO) Respiratory Guidelines. These documents were selected for their clarity, comprehensive coverage, and relevance to low-resource clinical settings.

From the MSF manual, only Chapter 2 (Respiratory diseases) was indexed for this study rather than the full manual. The WHO respiratory guideline was indexed in full (54 pages). In total, 53 MSF respiratory pages and 54 WHO pages were processed. After chunking and preprocessing, 310 guideline chunks were generated and embedded into the vector store.


\begin{table}[!ht]
\centering
\caption{Guideline Sources and Indexed Content Summary}
\label{tab:guideline_summary}
\renewcommand{\arraystretch}{1.2}
\makebox[\textwidth][c]{%
\resizebox{0.8\textwidth}{!}{%
\begin{tabular}{|l|p{8cm}|}
\hline
\textbf{Item} & \textbf{Details} \\
\hline
MSF Source & Clinical Guidelines -- Diagnosis and Treatment Manual (Respiratory Chapter Only) \\
\hline
WHO Source & Respiratory Guidelines -- Medicines Used in Respiratory Diseases \\
\hline
MSF Section Used & Chapter 2: Respiratory Diseases \\
\hline
WHO Section Used & Full Respiratory Guideline Document \\
\hline
MSF Pages Indexed & 53 pages \\
\hline
WHO Pages Indexed & 54 pages \\
\hline
Total Pages Processed & 107 pages \\
\hline
Total Guideline Chunks Generated & 310 chunks \\
\hline
\end{tabular}
}
}
\end{table}




While MSF and WHO guidelines provide internationally recognized clinical standards, they are designed primarily for general and resource-variable healthcare settings rather than institution-specific tertiary hospital workflows. They may therefore not fully reflect subspecialty-level diagnostic pathways or locally adapted treatment protocols. The evaluation in this study measures alignment with internationally standardized respiratory guidance rather than hospital-specific clinical practice.

All experiments were conducted using the implemented Python-based application and its integrated evaluation workflow. The complete system, including the retrieval modules, agentic orchestration logic, and web interface, was executed as a unified pipeline that accepted structured clinical inputs and produced end-to-end second-opinion outputs. Model inference was performed through hosted API services, while retrieval, orchestration, and evaluation logic were executed within the application environment.

Four representative clinical scenarios were designed to evaluate system behavior across different patient profiles. The evaluation cases were selected to reflect clinically complex situations in which second medical opinions are commonly sought, with emphasis on system behavior rather than clinical effectiveness. The full case descriptions are available in the public repository\footnote{\url{https://github.com/djh97/smo/blob/main/data/samples/cases.json}}:

\begin{itemize}
    \item A 6-year-old child with severe asthma exacerbation.
    \item A 1-year-old malnourished infant with post-measles pneumonia.
    \item A 9-year-old child with mild asthma symptoms.
    \item A 70-year-old adult with severe community-acquired pneumonia.
\end{itemize}

The inclusion of both pediatric and adult cases was intended to evaluate whether the retrieval and reasoning pipeline adapts to age-specific guideline criteria, severity thresholds, and dosing recommendations. Respiratory management protocols frequently differentiate treatment strategies across age groups, making mixed-age scenarios appropriate for assessing structured, guideline-aligned second-opinion generation. However, results remain limited to the respiratory domain and do not imply generalizability to unrelated clinical specialties without additional domain-specific validation.

All clinical scenarios were constructed as simulated cases to enable controlled evaluation of retrieval accuracy, orchestration behavior, and output structure without involving real patient records. This design isolates architectural performance from real-world clinical variability and aligns with early-stage artifact validation in design science research.

The evaluation design used two complementary testing procedures:

\begin{itemize}
    \item \textbf{RAG Functional Testing:} Each LLM was tested with and without retrieval augmentation to assess the influence of guideline-based evidence on output quality and alignment.
    
    \item \textbf{Agentic Output Testing:} The final synthesized output generated by the controller was compared with the RAG-enhanced responses from individual LLMs. Evaluation criteria included internal consistency, treatment completeness, and clinical appropriateness.
\end{itemize}




\begin{table*}[!ht]
\caption{Observed Clinical Differences Between Baseline and RAG Outputs in Case 1, a 6-Year-Old with Severe Asthma}
\label{tab:rag_clinical_enhancements}
\renewcommand{\arraystretch}{1.5}
\setlength{\tabcolsep}{4pt}
\normalsize
\resizebox{\textwidth}{!}{
\begin{tabular}{|>{\centering\arraybackslash}p{2.5cm}|c|p{5.3cm}|p{10.7cm}|}
\hline
\textbf{Model} & \textbf{Category} & \textbf{Baseline Output} & \textbf{RAG-Enhanced Improvement} \\
\hline

\multirow{7}{*}{\centering GPT-4o}
& Guideline Source & Not mentioned & WHO and MSF guidelines explicitly cited, increasing clinical credibility \\
\cline{2-4}
& Oxygen Target & Maintain $\mathrm{SpO_2}$ above 92\% & Recommends more conservative target: maintain $\mathrm{SpO_2}$ above 94\% based on WHO/MSF \\
\cline{2-4}
& Oxygen Delivery & Not specified & Specifies nasal cannula or face mask, adding clarity for implementation \\
\cline{2-4}
& ICU Admission & General escalation mention & Adds ICU admission conditions such as no improvement or signs of deterioration \\
\cline{2-4}
& Antibiotics & Not discussed & Introduces consideration of antibiotics if bacterial infection is suspected \\
\cline{2-4}
& Monitoring & General monitoring suggested & Advises structured monitoring: vital signs, treatment response, reassessment intervals \\
\cline{2-4}
& Follow-up & Review long-term plan & Recommends formal follow-up appointment post-stabilization \\
\hline

\multirow{6}{*}{\centering \makecell{Claude \\ Opus 4.6}}
& Guideline Source & Not mentioned & WHO and MSF criteria referenced for severity classification \\
\cline{2-4}
& Severity Thresholds & Based on symptoms ($\mathrm{SpO_2}$, HR) & Reinforces objective thresholds explicitly using guideline-based language \\
\cline{2-4}
& Bronchodilator Dosing & Every 20 mins for 1 hour & Matched in both, but RAG reinforces adherence to specific pediatric protocol \\
\cline{2-4}
& Corticosteroid Dosing & No dose or cap mentioned & Adds specific dose: 1--2 mg/kg/day with maximum limit \\
\cline{2-4}
& Admission Decision & Hospitalization suggested if no improvement & Specifies criteria for admission based on severity and $\mathrm{SpO_2}$ levels \\
\cline{2-4}
& Education and Discharge & Review action plan and meds & Adds education on inhaler use, discharge plan, and follow-up timing \\
\hline

\multirow{6}{*}{\centering \makecell{Gemini 2.5 \\ Flash}}
& Guideline Source & Not referenced & Adds structured, guideline-aligned treatment sections \\
\cline{2-4}
& Severity Classification & Classifies as moderate--severe & RAG confirms severity with stricter language based on thresholds ($\mathrm{SpO_2}$ $<$ 90\%) \\
\cline{2-4}
& Oxygen Therapy & Target $\mathrm{SpO_2}$ 94--98\% & Retained in RAG, reinforced with guideline reference \\
\cline{2-4}
& Medication Structure & Well-structured in baseline & RAG adds additional structure such as dosing frequency, delivery method, and action plan \\
\cline{2-4}
& ICU Escalation & Mentioned generally & Adds explicit ICU triggers: altered mental status, exhaustion, or mechanical ventilation need \\
\cline{2-4}
& Education and Follow-up & Includes plan review and referral & RAG improves formatting and structure; core content similar \\
\hline
\end{tabular}
}
\end{table*}

\subsection{RAG Functional Testing Across LLMs}

To examine the effect of retrieval grounding, we compared the outputs of GPT-4o, Claude Opus 4.6, and Gemini 2.5 Flash with and without RAG for a representative case involving a 6-year-old child with a severe asthma exacerbation. This case was used to observe how retrieved guideline evidence changed the content and structure of model responses.

A summary of the observed differences is presented in Table~\ref{tab:rag_clinical_enhancements} for Case~1. For GPT-4o, RAG introduced explicit oxygen targets, detailed delivery methods using nasal cannula or face mask, and ICU admission criteria based on non-response to initial treatment. Claude Opus 4.6 showed clearer bronchodilator dosing intervals, quantified corticosteroid recommendations, and structured discharge planning including follow-up visits. Gemini 2.5 Flash showed stronger severity framing based on oxygen saturation thresholds, more structured medication schedules, and escalation indicators such as signs of respiratory fatigue or altered mental status.

These observations show how retrieval grounding can make model responses more specific, more structured, and more closely aligned with WHO and MSF guidance in a clinically meaningful example.

%The full set of raw responses, including both baseline and RAG-enhanced outputs for all four cases, has been uploaded to the project's GitHub repository. This ensures transparency and enables reproducibility. The remaining three cases were evaluated using the same criteria. Across all scenarios, RAG consistently improved clarity, clinical validity, and structure of the LLM outputs.



\subsection{Agentic Pipeline Testing and Output Review}
Agentic pipeline evaluation examines the system's ability to produce consistent, guideline-aligned outputs across varied clinical scenarios, including both initial and follow-up encounters. This evaluation focuses on system behavior and reliability and is not intended to assess clinical efficacy or patient outcomes.

Across all four test cases, the agentic pipeline accepted structured inputs, retrieved guideline context, invoked the selected expert tools, synthesized a final response, and rendered the output in the interface. Both new and follow-up scenarios completed successfully, which indicates coherent end-to-end operation across the evaluated respiratory cases.

\begin{figure*}[!ht]
\centering
\setlength{\fboxsep}{0pt} 
\setlength{\fboxrule}{0.4pt} 
\fbox{
\includegraphics[width=\textwidth]{UI_NewVisit.jpg}}
\caption{Final agentic opinion synthesized for new visit case}
\label{fig:case1_output}
\end{figure*} 

Figure~\ref{fig:case1_output} shows the output of the system for Case 1. This case involved a 6-year-old child with a severe asthma exacerbation. The figure displays both the input interface and the final structured output generated by the agentic system.

The agentic controller correctly recognized this as a new case by detecting the absence of a previously stored patient ID and the visit type marked as ``New''. Based on this context, the system initiated a fresh diagnostic assessment using its retrieval-augmented expert tools. Depending on the case, the controller may consult one, two, or all available tools before synthesis.

In this example, the controller generated intermediate tool outputs conditioned on the retrieved WHO/MSF guideline context and then synthesized a unified structured opinion. The final response (logged as the controller step) integrates diagnosis, severity classification, and a treatment plan into the standardized four-part format displayed in the interface. This confirms that the system supports dynamic tool routing while maintaining a single, clinician-facing output.


\begin{figure*}[!ht]
\centering
\setlength{\fboxsep}{0pt} 
\setlength{\fboxrule}{0.4pt} 
\fbox{
\includegraphics[width=\textwidth]{UI_Followup.jpg}}
\caption{Final agentic opinion synthesized for follow-up visit}
\label{fig:followup-case}
\end{figure*}

The agentic system was also tested with a follow-up scenario involving the same patient. As shown in Figure~\ref{fig:followup-case}, the system automatically detected the existing patient ID and correctly identified the visit as a follow-up. It then compared the new vital signs with prior values and assessed progression using the prior visit context available through the application's local stored visit record. In this case, the system identified clinical deterioration. The output classified the condition as ``worsening'' and generated an updated treatment plan. This confirmed the agentic controller's ability to support continuity across visits within the released research prototype.

For transparency, the representative case was logged during evaluation in a four-case execution trace that preserves the retrieved guideline context, prompts, intermediate tool outputs, controller decisions, and the final synthesized response. This trace is provided in the Supplementary Material in the public repository.

\section{Results and Discussion}
\label{sec:discussion}
This section presents the evaluation results for our agentic AI system and discusses its performance in producing evidence-based second medical opinions. We include guideline alignment, classification metrics, cost efficiency, latency, repeatability, and comparisons with other systems. It also considers practical feasibility, clinical integration, safety, and scalability based on observed system behavior.

\subsection{Guideline Alignment Performance}

Alignment with established clinical guidelines is a central requirement for safe and trustworthy decision support systems. In the context of automated second medical opinion generation, consistency with authoritative WHO/MSF recommendations provides an operational proxy for clinical soundness and reasoning adequacy. This subsection evaluates alignment from two complementary perspectives: semantic similarity to guideline-derived reference answers and reasoning-level agreement under G-Eval. Within the semantic analysis, per-case results, aggregate summary statistics, and visual comparisons are presented to capture both average performance and cross-case variability across configurations.

\subsubsection{Semantic alignment using cosine similarity}

To quantify this alignment, we computed cosine similarity between model-generated recommendations and structured WHO/MSF-aligned reference answers using the \texttt{all-MiniLM-L6-v2} sentence-transformer model.

We evaluated three configurations across the predefined four clinically representative respiratory cases in Section \ref{sec:test}: (i) single-model baseline without RAG, (ii) single-model RAG, and (iii) the proposed multi-LLM agentic pipeline that synthesizes outputs from GPT-4o, Claude Opus 4.6, and Gemini 2.5 Flash. For each configuration, alignment was measured against a fixed guideline-based reference answer for the corresponding case.

Table~\ref{tab:alignment_scores_structured} reports per-case cosine alignment scores. These values represent semantic proximity between generated recommendations and WHO/MSF-derived reference standards, and they provide a case-level view of how alignment changes across baseline, RAG, and agentic configurations.

\begin{table*}[!ht]
\centering
\caption{Cosine alignment between model outputs and WHO/MSF reference answers (4 Clinical Cases)}
\label{tab:alignment_scores_structured}
\renewcommand{\arraystretch}{1.2}
\small
\begin{tabular}{|c|cc|cc|cc|c|}
\hline
\textbf{Case} 
& \multicolumn{2}{c|}{\textbf{GPT-4o}} 
& \multicolumn{2}{c|}{\textbf{Claude Opus 4.6}} 
& \multicolumn{2}{c|}{\textbf{Gemini 2.5 Flash}} 
& \textbf{Agentic} \\
\cline{2-7}
& Baseline & RAG & Baseline & RAG & Baseline & RAG &  \\
\hline
P001 & 0.7223 & 0.7531 & 0.6442 & 0.6902 & 0.6988 & 0.6864 & 0.8011 \\
P002 & 0.7648 & 0.6549 & 0.6207 & 0.5903 & 0.7701 & 0.7191 & 0.7583 \\
P003 & 0.7264 & 0.8131 & 0.7706 & 0.7638 & 0.8035 & 0.7846 & 0.8031 \\
P004 & 0.7288 & 0.6653 & 0.5516 & 0.6187 & 0.6067 & 0.6174 & 0.6892 \\
\hline
\end{tabular}
\end{table*}

To examine overall performance trends, we computed mean alignment and cross-case variability for each configuration. These results are summarized numerically in Table~\ref{tab:alignment_summary_structured}, which complements the per-case view by condensing performance into aggregate statistics.

\begin{table}[!ht]
\centering
\caption{Mean cosine alignment and cross-case variability}
\label{tab:alignment_summary_structured}
\renewcommand{\arraystretch}{1.2}
\small
\begin{tabular}{|l|c|c|}
\hline
\textbf{Configuration} & \textbf{Mean} & \textbf{Std} \\
\hline
Agentic Pipeline & 0.7629 & 0.0533 \\
GPT-4o Baseline & 0.7356 & 0.0197 \\
GPT-4o RAG & 0.7216 & 0.0753 \\
Claude Opus 4.6 Baseline & 0.6468 & 0.0914 \\
Claude Opus 4.6 RAG & 0.6658 & 0.0777 \\
Gemini 2.5 Flash Baseline & 0.7198 & 0.0871 \\
Gemini 2.5 Flash RAG & 0.7019 & 0.0696 \\
\hline
\end{tabular}
\end{table}

Figure~\ref{fig:mean_alignment} visualizes mean alignment together with standard deviation error bars, while Figure~\ref{fig:alignment_variability} shows cross-case dispersion across configurations. Together with the tabulated results, these visualizations help illustrate both overall performance trends and variability across cases.

\begin{figure}[!ht]
\centering
\includegraphics[width=\textwidth]{alignment_std.pdf}
\caption{Mean cosine alignment (\(\pm\) standard deviation) across four predefined respiratory cases under baseline, RAG, and agentic configurations.}
\label{fig:mean_alignment}
\end{figure}

\begin{figure}[!ht]
\centering
\includegraphics[width=\textwidth]{alignment_variability.pdf}
\caption{Cross-case alignment variability (standard deviation) across four predefined respiratory cases under baseline, RAG, and agentic configurations.}
\label{fig:alignment_variability}
\end{figure}


Across the four cases, retrieval augmentation did not uniformly improve cosine alignment relative to baseline performance. GPT-4o shows a lower mean alignment under RAG (0.7216) than baseline (0.7356), and Gemini 2.5 Flash also decreases under RAG (0.7019) relative to baseline (0.7198). In contrast, Claude Opus 4.6 shows a modest improvement under RAG (0.6658) compared with baseline (0.6468), although both settings exhibit higher dispersion than GPT-4o baseline.

The agentic pipeline achieves the highest mean cosine alignment among the evaluated configurations (0.7629) while maintaining moderate cross-case variability (Std 0.0533). At the case level, the agentic configuration achieves the strongest alignment on P001, remains competitive on P002 and P003, and outperforms both Claude configurations and Gemini RAG on P004, although it is not the top-scoring configuration in every individual case. Overall, these results indicate that controller-based multi-LLM orchestration can preserve strong semantic agreement with guideline-derived reference answers and deliver more favorable average alignment than the individual baseline and single-model RAG configurations in the evaluated respiratory scenarios.


\subsubsection{G-Eval reasoning-level agreement analysis}

While cosine similarity quantifies semantic proximity to guideline-derived reference answers, clinical validity also requires evaluation of reasoning structure, management completeness, and the ability to reject clinically mismatched references. To assess this dimension, we employed the G-Eval framework, which uses a large language model as a structured evaluator. Each generated output was assigned a continuous agreement score in the interval [0,1].

The G-Eval score \( G \) is defined as:

\begin{equation}
G = \text{LLM}_{\text{judge}}\left(\text{input}, \text{actual\_output}, \text{expected\_output}\right)
\label{eq:geval_score}
\end{equation}

GPT-4o was used as the judging model based on prior evidence demonstrating its structured reasoning capability in expert review settings \cite{r2,r3}. The evaluation prompt required the judge to assess diagnostic correctness, severity classification, treatment completeness, patient-specificity, and escalation appropriateness relative to WHO/MSF standards.

Each model output was evaluated against three reference pair types for each clinical case: (i) the correct WHO/MSF-aligned reference answer (positive pair), (ii) a clinically related but incorrect reference from the same broader respiratory family (near-negative pair), and (iii) a clearly mismatched reference from a different condition family (far-negative pair). Binary judgments were then obtained by thresholding the continuous G-Eval score. The operating threshold was set to \(\tau = 0.90\), corresponding to the ``strong agreement'' band in the evaluation rubric. This threshold was selected to prioritize high-confidence acceptance in a safety-sensitive clinical context, consistent with recommendations for strict decision boundaries in medical AI systems \cite{r4}. Under this criterion, only outputs with minor omissions or harmless wording differences were classified as agreement.

\begin{equation}
\text{Judgment} =
\begin{cases}
    \text{Agree} & \text{if } G \geq \tau \\
    \text{Diverge} & \text{otherwise}
\end{cases}
\label{eq:threshold_judgment}
\end{equation}

Figure~\ref{fig:threshold_curve} shows how true positives, false positives, true negatives, and false negatives vary across thresholds over the complete set of positive and negative evaluation pairs. The threshold sweep indicates a marked tradeoff in the 0.85--0.90 region, with \(\tau = 0.90\) preserving complete rejection of negative pairs while retaining the most strongly aligned positive pairs.

\begin{figure}[!ht]
    \centering
    \includegraphics[width=\linewidth]{threshold.png}
    \caption{Confusion-matrix elements across varying G-Eval thresholds over positive, near-negative, and far-negative evaluation pairs.}
    \label{fig:threshold_curve}
\end{figure}

Under this definition:

\begin{itemize}
    \item True positive (TP): a positive pair classified as Agree.
    \item False negative (FN): a positive pair classified as Diverge.
    \item False positive (FP): a negative pair classified as Agree.
    \item True negative (TN): a negative pair classified as Diverge.
\end{itemize}

The standard classification metrics are defined as follows:

\begin{equation}
\text{Precision} = \frac{TP}{TP + FP}
\end{equation}

\begin{equation}
\text{Recall} = \frac{TP}{TP + FN}
\end{equation}

\begin{equation}
\text{Specificity} = \frac{TN}{TN + FP}
\end{equation}

\begin{equation}
\text{F1-score} = 2 \times \frac{\text{Precision} \times \text{Recall}}{\text{Precision} + \text{Recall}}
\end{equation}

Figure~\ref{fig:confusion_matrices} presents confusion matrices for all evaluated configurations under the selected threshold \(\tau = 0.90\). These matrices illustrate how often each system accepts strongly aligned reference pairs while rejecting mismatched ones.

\begin{figure*}[!ht]
    \centering
    \subfloat[GPT-4o]{%
        \includegraphics[width=0.45\linewidth]{gptmatrix.png}}%
    \hfill
    \subfloat[Claude Opus 4.6]{%
        \includegraphics[width=0.45\linewidth]{claudematrix.png}}\\[1ex]
    \subfloat[Gemini 2.5 Flash]{%
        \includegraphics[width=0.45\linewidth]{geminimatrix.png}}%
    \hfill
    \subfloat[Agentic]{%
        \includegraphics[width=0.45\linewidth]{agenticmatrix.png}}%
    \caption{Confusion matrices under the selected G-Eval threshold \(\tau = 0.90\) for GPT-4o, Claude Opus 4.6, Gemini 2.5 Flash, and the agentic pipeline.}
    \label{fig:confusion_matrices}
\end{figure*}

Positive-pair case-level G-Eval scores are reported in Table~\ref{tab:geval_all_cases}. Claude Opus 4.6 and the agentic pipeline achieve the highest mean positive-pair score (0.925 each). The agentic pipeline matches Claude on severe asthma, severe pneumonia, and elderly CAP (0.95 in each case), while Gemini 2.5 Flash attains the highest score on mild asthma (0.95). GPT-4o and Gemini 2.5 Flash share the lowest mean positive-pair score (0.775), with GPT-4o scoring lowest on severe asthma and elderly CAP.

\begin{table*}[!ht]
\caption{Positive-pair case-level G-Eval scores}
\label{tab:geval_all_cases}
\centering
\renewcommand{\arraystretch}{1.25}
\setlength{\tabcolsep}{8pt}
\small
\resizebox{\textwidth}{!}{
\begin{tabular}{|l|c|c|c|c|}
\hline
\textbf{Clinical Case} & \textbf{GPT-4o} & \textbf{Claude Opus 4.6} & \textbf{Gemini 2.5 Flash} & \textbf{Agentic} \\
\hline
\textbf{Severe asthma} & 0.70 & 0.95 & 0.70 & 0.95 \\
\hline
\textbf{Severe pneumonia} & 0.85 & 0.95 & 0.70 & 0.95 \\
\hline
\textbf{Mild asthma} & 0.85 & 0.85 & 0.95 & 0.85 \\
\hline
\textbf{Elderly CAP} & 0.70 & 0.95 & 0.75 & 0.95 \\
\hline
\end{tabular}}
\end{table*}

The quantitative implications of the thresholded G-Eval results are reported in Table~\ref{tab:metrics_llms}. Under the selected threshold, Claude Opus 4.6 and the agentic pipeline each achieve precision \(=1.00\), recall \(=0.75\), specificity \(=1.00\), and F1-score \(=0.86\). Gemini 2.5 Flash also rejects all negative pairs and retains one positive pair above the 0.90 cutoff, yielding precision \(=1.00\), recall \(=0.25\), specificity \(=1.00\), and F1-score \(=0.40\). GPT-4o rejects all negative pairs, but none of its positive-pair scores reach the 0.90 high-agreement cutoff.

\begin{table}[!ht]
\centering
\caption{Thresholded G-Eval metrics under the selected threshold \(\tau = 0.90\)}
\label{tab:metrics_llms}
\renewcommand{\arraystretch}{1.2}
\small
\begin{tabular}{|l|c|c|c|c|}
\hline
\textbf{Model} & \textbf{Precision} & \textbf{Recall} & \textbf{Specificity} & \textbf{F1-score} \\
\hline
GPT-4o & 0.00 & 0.00 & 1.00 & 0.00 \\
Claude Opus 4.6 & 1.00 & 0.75 & 1.00 & 0.86 \\
Gemini 2.5 Flash & 1.00 & 0.25 & 1.00 & 0.40 \\
Agentic & 1.00 & 0.75 & 1.00 & 0.86 \\
\hline
\end{tabular}
\end{table}

These thresholded results show that the agentic pipeline remains among the strongest systems under a stringent high-agreement acceptance rule, matching Claude Opus 4.6 on both mean positive-pair agreement and thresholded metrics. This indicates that multi-LLM orchestration can preserve high-quality reasoning while enforcing strict rejection of mismatched reference patterns.

These results complement the cosine similarity analysis. Whereas cosine similarity captures semantic proximity to the reference text, G-Eval evaluates whether the reasoning structure and management plan meet a predefined standard of completeness. The combined findings therefore support the conclusion that multi-LLM orchestration improves reasoning-level agreement with guideline standards while maintaining stability across heterogeneous clinical scenarios.

\subsubsection{Error analysis and handling of misaligned recommendations}

The cosine similarity and G-Eval analyses quantify how closely system outputs follow WHO/MSF guideline-based reference answers, but they do not guarantee absence of errors. In a safety-critical setting such as clinical decision support, it is important to examine how misalignments arise and to clarify how the system is intended to be used when they occur.

A review of the pair-level G-Eval results suggests two recurrent misalignment patterns for weaker configurations. The first is under-specification, where the diagnosis and severity classification are broadly correct but important management details are incomplete, such as escalation criteria, antibiotic specificity, oxygen targets, or follow-up instructions. The second is same-family over-generalization, where a response remains plausible at the respiratory-disease level but does not adequately preserve patient-specific features such as age group, comorbidity context, or case-specific treatment pathway when compared against a near-negative reference. Table~\ref{tab:error_patterns} summarizes these patterns.

\begin{table}[!ht]
\centering
\caption{Observed error patterns in misaligned recommendations}
\label{tab:error_patterns}
\renewcommand{\arraystretch}{1.2}
\small
\resizebox{\textwidth}{!}{
\begin{tabular}{|>{\centering\arraybackslash}p{3.2cm}|p{5.2cm}|p{5.6cm}|}
\hline
\textbf{Pattern} & \textbf{Description} & \textbf{Observed implication} \\
\hline
Under-specification & The response is broadly correct at the diagnosis or severity level, but omits clinically important management details. & Outputs may appear semantically aligned while still lacking treatment-critical content such as escalation criteria, monitoring detail, oxygen targets, or follow-up instructions. \\
\hline
Same-family over-generalization & The response stays within the correct respiratory disease family but does not preserve the patient-specific management context. & Outputs may sound clinically plausible yet map to the wrong age group, severity context, comorbidity profile, or treatment pathway when compared with a near-negative reference. \\
\hline
\end{tabular}
}
\end{table}

These patterns help explain the difference between cosine similarity and G-Eval. Semantic alignment can remain high because core diagnostic content overlaps with the reference answer, while G-Eval marks the same response as weaker if it omits treatment-critical details or fails to preserve patient-specific management logic. Under the selected threshold and reported evaluation setting, all evaluated systems reject the negative pairs, but only Claude Opus 4.6 and the agentic pipeline retain three of the four positive pairs within the strong-agreement region.

The agentic configuration appears to reduce these omission and over-generalization patterns by synthesizing across dynamically selected expert-tool outputs while preserving retrieved guideline context in the final controller step. In the present evaluation, this yields the joint-highest mean positive-pair score and matches the strongest single-model configuration under the stringent 0.90 threshold, indicating that controller-based synthesis can achieve strong reasoning-level agreement without increasing false acceptance of mismatched references.

Despite these improvements, the system is not intended to operate as an autonomous diagnostic authority. It is designed as a decision-support tool that presents structured, guideline-grounded recommendations in a read-only format. Final responsibility for diagnosis, severity assessment, and treatment decisions remains with the attending clinician. When AI-generated recommendations conflict with local protocols or clinician judgment, the clinician's assessment and institutional guidelines must take precedence. Safe deployment would also benefit from structured audit logging, periodic review of recurring error patterns, and clinician-in-the-loop studies that evaluate how practitioners interpret and act on AI suggestions in real clinical workflows.

\subsection{Operational performance analysis}
\label{sec:operational}

Beyond alignment with clinical guidelines, practical deployment of an agentic decision-support system requires evaluation of its computational efficiency, latency, and output stability. This subsection examines the operational characteristics of the proposed framework, including token-level cost, processing time, and repeatability under repeated executions.


\subsubsection{Computational Cost}

To evaluate the economic feasibility of the agentic pipeline for generating second medical opinions, we measured token-based inference costs for four representative clinical cases (P001--P004). The pipeline relies on three models: GPT-4o\footnote{\url{https://openai.com/api/pricing}}, Claude Opus 4.6\footnote{\url{https://www.anthropic.com/pricing}}, and Gemini 2.5 Flash\footnote{\url{https://ai.google.dev/gemini-api/docs/pricing}}.

Token and cost logging were captured from a final instrumented execution. GPT-4o token counts were computed exactly using \texttt{tiktoken}, an OpenAI tokenizer library that provides model-aware token counting. Claude Opus 4.6 and Gemini 2.5 Flash token counts were estimated using \texttt{cl100k\_base}, a byte-pair encoding distributed with \texttt{tiktoken} that was used here as a common approximation across non-OpenAI providers; these estimated counts are reflected in Table~\ref{tab:agentic_cost_final}.

Each case execution consists of a dynamic sequence of retrieval-augmented tool calls followed by a final controller synthesis step. The controller may invoke one, two, or all three expert tools depending on the case, and then returns a single consolidated structured recommendation via a final GPT-4o controller response (logged as \texttt{role="controller"}). Consequently, per-case cost depends on which tools are used rather than assuming a fixed three-tool path.

Cost calculations use published API pricing as of mid-2025:

\begin{itemize}
    \item \textbf{GPT-4o}: USD~2.50 per 1 million input tokens and USD~10.00 per 1 million output tokens.
    \item \textbf{Claude Opus 4.6}: USD~5.00 per 1 million input tokens and USD~25.00 per 1 million output tokens.
    \item \textbf{Gemini 2.5 Flash}: USD~0.30 per 1 million input tokens and USD~2.50 per 1 million output tokens.
\end{itemize}

The cost for each model is calculated using:

\begin{equation}
\text{Cost}_{\text{LLM}} = \left( \frac{T_{\text{in}}}{10^6} \times P_{\text{in}} \right) + \left( \frac{T_{\text{out}}}{10^6} \times P_{\text{out}} \right)
\label{eq:llm_cost}
\end{equation}

\noindent
Here, \( T_{\text{in}} \) and \( T_{\text{out}} \) denote the number of input and output tokens respectively. The parameters \( P_{\text{in}} \) and \( P_{\text{out}} \) represent the token prices in USD.

\noindent
The total cost of the agentic pipeline for case \(c\) is defined as:

\begin{equation}
\text{Cost}_{\text{Agentic}}(c) = \sum_{m \in \mathcal{S}(c)} \text{Cost}_{m}(c) + \text{Cost}_{\text{Controller}}(c)
\label{eq:agentic_cost}
\end{equation}

\noindent
where $\mathcal{S}(c)$ denotes the subset of expert tools invoked for clinical case $c$ (one to three tools), and the controller term represents the final GPT-4o synthesis call that produces the unified structured recommendation.

The size of the underlying guideline corpus does not directly determine per-case cost because retrieval is restricted to a fixed number of top-ranked chunks (top-$k=3$). However, longer retrieved passages increase prompt length and therefore token consumption during tool invocation.

Table~\ref{tab:agentic_cost_final} presents the detailed token usage and cost breakdown per case from the final instrumented execution record used for analysis. This record includes only agentic pipeline calls (dynamic tool invocations plus the final controller synthesis step) and explicitly excludes baseline and single-model RAG runs. The table reports the tools used per case and shows that the controller can skip tools when sufficient evidence is obtained (e.g., P003 omits the Gemini tool), which reduces total cost relative to a fixed three-tool assumption.

\begin{table*}[!ht]
\caption{Input and output columns report tokens, and cost columns report USD per case under dynamic tool invocation. GPT-4o token counts are computed exactly, while Claude Opus 4.6 and Gemini 2.5 Flash token counts are estimated.}
\label{tab:agentic_cost_final}
\centering
\renewcommand{\arraystretch}{1.2}
\footnotesize
\resizebox{\textwidth}{!}{
\begin{tabular}{|c|p{2.6cm}|ccc|ccc|ccc|ccc|c|}
\hline
\textbf{Case} 
& \textbf{Tools Used}
& \multicolumn{3}{c|}{\textbf{GPT-4o Tool}} 
& \multicolumn{3}{c|}{\textbf{Claude Tool (est)}} 
& \multicolumn{3}{c|}{\textbf{Gemini Tool (est)}} 
& \multicolumn{3}{c|}{\textbf{Controller (GPT-4o)}} 
& \textbf{Total} \\
\hline
&
& \textbf{In} & \textbf{Out} & \textbf{Cost} 
& \textbf{In} & \textbf{Out} & \textbf{Cost} 
& \textbf{In} & \textbf{Out} & \textbf{Cost} 
& \textbf{In} & \textbf{Out} & \textbf{Cost} 
& \textbf{(USD)} \\
\hline
P001 & \makecell[c]{[All]}
& 594 & 311 & 0.004595 & 607 & 1009 & 0.028260 & 607 & 351 & 0.001060 & 2916 & 534 & 0.012630 & \textbf{0.046545} \\
P002 & \makecell[c]{[All]}
& 612 & 264 & 0.004170 & 622 & 1132 & 0.031410 & 622 & 363 & 0.001094 & 3023 & 683 & 0.014388 & \textbf{0.051062} \\
P003 & \makecell[c]{[GPT + Claude]}
& 591 & 318 & 0.004658 & 604 & 942 & 0.026570 & 0 & 0 & 0.000000 & 2500 & 359 & 0.009840 & \textbf{0.041068} \\
P004 & \makecell[c]{[All]}
& 574 & 286 & 0.004295 & 585 & 1126 & 0.031075 & 585 & 546 & 0.001541 & 3185 & 576 & 0.013722 & \textbf{0.050633} \\
\hline
\end{tabular}
}
\end{table*}





Figure~\ref{fig:agentic_cost_bar} summarizes the average cost contribution of each agentic component across the four cases. In this run, the average component costs were USD~0.0044 for the GPT-4o tool call, USD~0.0293 for the Claude tool call, USD~0.0009 for the Gemini tool call, and USD~0.0126 for the controller synthesis, yielding an average full-pipeline cost of USD~0.0473.

\begin{figure}[!ht]
\centering
\includegraphics[width=\textwidth]{agentic_cost_avg_dynamic.png}
\caption{Average token-based cost per component across four cases, reflecting dynamic tool invocation (one to three expert tools) plus the final GPT-4o controller stages.}
\label{fig:agentic_cost_bar}
\end{figure}

Extrapolation to 1,000 patient consultations per day gives an estimated daily cost of about USD~47.33, based on the observed average full-pipeline cost of USD~0.0473 per case. In practice, whether this operational cost is acceptable depends on local budget constraints and the expected clinical and operational value of decision support; such considerations are commonly addressed through economic evaluation frameworks for clinical decision support systems \cite{White2023}. The estimate remains sensitive to tool selection because dynamic agent execution can skip certain tools when sufficient evidence is obtained earlier, which reduces overall cost.


\subsubsection{Latency and Time Analysis}

The total processing time for generating an agentic medical opinion includes multiple sequential steps. These steps reflect the core stages of the agentic pipeline: retrieval of clinical context, reasoning across multiple LLMs, and synthesis of a final output. 

The total processing time per case is modeled using Equation~\ref{eq:time_total}, which sums the duration of each core stage:

\begin{equation}
T_{\text{Total}} = T_{\text{Retrieval}} + \sum_{i=1}^{n} T_{\text{LLM}_i} + T_{\text{Controller}}
\label{eq:time_total}
\end{equation}

Here, \( T_{\text{Retrieval}} \) denotes the cumulative time spent retrieving guideline context across the controller and any invoked expert tools, \( T_{\text{LLM}_i} \) is the response time for each expert LLM used (where \( n \leq 3 \)), and \( T_{\text{Controller}} \) is the cumulative controller time, including routing and final synthesis. Since the controller may use one, two, or all three LLMs based on the required level of reasoning, the realized runtime depends on the dynamic path selected for each case. 

Table~\ref{tab:agentic_time} summarizes the time breakdown across four clinical cases from the instrumented run log. Each row provides the cumulative retrieval time, per-model reasoning time for the tools that were actually invoked, controller time, and the total end-to-end runtime. Figure~\ref{fig:agentic_time_chart} shows the average time contribution of each component in the agentic pipeline. In this run, the largest share of processing time comes from Claude's reasoning duration, followed by the controller stages, while cumulative retrieval remains comparatively small.


\begin{table*}[!ht]
\caption{Processing time breakdown for the agentic pipeline per case under dynamic tool invocation}
\label{tab:agentic_time}
\centering
\renewcommand{\arraystretch}{1.2}
\footnotesize
\resizebox{\textwidth}{!}{%
\begin{tabular}{|c|p{2.6cm}|c|c|c|c|c|c|}
\hline
\textbf{Case} & \textbf{Tools Used} & \textbf{Guideline Retrieval (s)} & \textbf{GPT-4o Tool (s)} & \textbf{Claude Tool (s)} & \textbf{Gemini Tool (s)} & \textbf{Controller (s)} & \textbf{Total Time (s)} \\
\hline
P001 & \makecell[c]{[All]} & 1.03 & 3.93 & 24.63 & 9.54 & 6.27 & 45.40 \\
P002 & \makecell[c]{[All]} & 1.04 & 2.82 & 27.39 & 6.72 & 7.53 & 45.49 \\
P003 & \makecell[c]{[GPT + Claude]} & 0.87 & 3.99 & 23.44 & 0.00 & 6.63 & 34.92 \\
P004 & \makecell[c]{[All]} & 1.75 & 6.30 & 28.60 & 8.38 & 7.81 & 52.84 \\
\hline
\textbf{Average} &  & \textbf{1.17} & \textbf{4.26} & \textbf{26.02} & \textbf{6.16} & \textbf{7.06} & \textbf{44.66} \\
\hline
\end{tabular}%
}
\end{table*}


\begin{figure}[!ht]
\centering
\includegraphics[width=\textwidth]{av_time.png}
\caption{Average latency per component across four cases under dynamic tool invocation}
\label{fig:agentic_time_chart}
\end{figure}

Unlike the earlier fixed-path approximation, this timing analysis is derived directly from instrumented runtime logs and therefore reflects the actual dynamic execution path selected per case. For example, P003 skips the Gemini tool entirely and records zero Gemini latency, which shortens its total runtime relative to the three-tool cases. Under this evaluated run, end-to-end latency ranges from 34.92 to 52.84 seconds, with a mean of 44.66 seconds across the four cases.

\subsubsection{Repeatability and Output Stability}
\label{sec:repeatability}

In addition to cross-case robustness, we evaluated intra-case stability by executing the full agentic pipeline 30 times on an identical moderate-to-severe asthma scenario. Each run used identical structured inputs and was evaluated against a fixed WHO/MSF-aligned reference answer.

Cosine similarity between generated outputs and the reference yielded a mean alignment of 0.8438 (SD = 0.0234), corresponding to a coefficient of variation of 2.8\%. Alignment scores ranged from 0.7791 to 0.8809, with a 95\% confidence interval for the mean of [0.8354, 0.8521]. Figure~\ref{fig:repeatability_all} summarizes repeatability using both a distributional view (boxplot) and a run-ordered view (run-to-run stability).

\begin{figure}[!ht]
\centering
\includegraphics[width=\textwidth]{repeatability_all.pdf}
\caption{Repeatability across 30 executions on an identical clinical case: distribution of cosine alignment scores (boxplot) and run-to-run alignment values, illustrating limited dispersion and no systematic drift under identical inputs.}
\label{fig:repeatability_all}
\end{figure}

The narrow interquartile range and absence of extreme outliers indicate limited semantic dispersion across repeated executions. The run-ordered plot further indicates no systematic drift across the 30 executions under identical inputs. Although minor lexical variation occurred between runs, diagnostic classification and severity assessment remained consistent.

These results demonstrate that the agentic framework produces reproducible and semantically stable guideline-aligned recommendations, a critical requirement for safety-sensitive clinical decision support systems.



\subsection{Practical Implications}

This section describes how the proposed agentic framework can move from theoretical design to real-world clinical practice. The practical implications focus on its application in clinical workflows and its potential to improve decision support and patient outcomes.

\subsubsection{Workflow Integration}

To support real-world deployment, our system is designed for seamless integration with existing clinical workflows. CDSS can adopt the agentic architecture to enhance clinician decision-making through LLM-powered reasoning grounded in trusted medical guidelines. The framework respects existing EHR infrastructures and can be invoked at the point of care to deliver structured, AI-generated second opinions.

This enables clinicians to validate or refine diagnoses and treatment plans with consistent and explainable outputs. Recent work in patient-centered LLM applications has emphasized the importance of contextual augmentation and structured reasoning in supporting medical professionals \cite{Neupane2025, MACIA2024147}. Our approach complements this direction by grounding recommendations in curated guideline content and generating structured outputs that support cautious clinical review.


\subsubsection{Improved Safety and Reliability}

Implementing this system aims to improve reliability by reducing the likelihood of unsupported recommendations through guideline-grounded retrieval and structured synthesis. The agentic pipeline queries multiple RAG-enabled LLM tools and uses a controller model to synthesize a unified recommendation from their outputs. While this design does not implement a runtime threshold-based acceptance gate, it supports safer decision support by providing structured and evidence-grounded outputs intended for clinician review rather than autonomous action.

\subsubsection{Clinical Decision Support Enhancement}

The framework delivers results that both explain the process and follow established guidelines instead of replacing human decision-making. The system delivers standardized evidence-based recommendations to clinicians which decreases their mental workload and practice variability but preserves their ability to make decisions. Research shows that effective decision support systems lead to better patient results and improved guideline compliance \cite{r34, XU2023130}. From a clinician trust perspective, the observed guideline alignment, evidence-grounded reasoning, and conservative thresholding are necessary but not sufficient; prospective usability studies, clinician-in-the-loop evaluations, and institutional governance processes would be required before routine deployment to confirm that the system's explanations, failure modes, and integration with local protocols are acceptable in practice.

\subsubsection{Scalability and Adaptability}

The framework features a modular structure which allows future modifications to include new therapeutic domains and updated recommendations. The approach maintains its relevance through healthcare knowledge advancement and provides enduring advantages. The system developers can modify retrieval sources and underlying models to accommodate new illnesses and therapies without rebuilding the entire system. The system maintains its operational capability across different clinical environments.


\subsection{Comparative Analysis of Existing Solutions}

\begin{table*}[!ht]
\caption{Technical and Evaluation Comparison of Related Work and Our Agentic AI Pipeline}
\label{tab:comparison_full}
\centering
\renewcommand{\arraystretch}{1.5}
\setlength{\tabcolsep}{4pt}
\normalsize
\resizebox{\textwidth}{!}{%
\begin{tabular}{|>{\centering\arraybackslash}p{2.5cm}|
                >{\centering\arraybackslash}p{2.7cm}|
                >{\centering\arraybackslash}p{2.8cm}|
                >{\centering\arraybackslash}p{3.8cm}|
                >{\centering\arraybackslash}p{3.8cm}|
                >{\centering\arraybackslash}p{2.7cm}|
                >{\centering\arraybackslash}p{4cm}|
                >{\centering\arraybackslash}p{2cm}|}
\hline
\textbf{System} & \textbf{Guideline Integration} & \textbf{Multi-Condition Support} & \textbf{Reasoning Approach} & \textbf{Architecture / Orchestration} & \textbf{External Knowledge / Retrieval} & \textbf{Evaluation Metrics} & \textbf{UI Implemented} \\
\hline

\cite{r10} & None & No (single domain) & NEPAR + ML classification & Static ML pipeline & No retrieval & Diagnostic accuracy evaluation & Yes \\
\hline
\cite{r11} & Limited, single-center & No (specific surgical decision) & Predictive modeling with patient-reported data & Single-task risk prediction & No retrieval & Protocol design; planned cost-utility analysis & Yes \\
\hline
\cite{r12} & None & No (binary classification) & Dense neural networks on tabular data & Single-task NN model & No retrieval & Benchmark dataset evaluation & No \\
\hline
\cite{r14} & No dynamic integration & Narrow prescription screening & XAI explanations (Local, Global, Counterfactual) & Single model with UI & No retrieval & Human-AI team performance study & Yes \\
\hline
\cite{r16} & Static external sources & Single-turn, single domain & RAG-enabled retrieval with LLM generation & Retriever and agent tools for multiple sources & Yes (PDFs, PubMed, CSV) & User satisfaction testing & Yes \\
\hline
\cite{r20} & No real-time guideline integration & Clinical trial-focused & Multi-agent with ReAct + Least-to-Most & Multi-agent GPT-4 system & No retrieval & Predictive performance (PR-AUC) evaluation & Yes \\
\hline
\cite{r22} & Static knowledge base & Multi-condition (case study) & Staged multi-agent reasoning & Modular agent framework & No dynamic retrieval & Case study demonstration & Yes \\
\hline
\textbf{Our Agentic AI Pipeline} & Dynamic retrieval of clinical guidelines & Supports new and follow-up visits & Retrieval-augmented, multi-LLM reasoning with structured synthesis & Agentic pipeline with orchestrated multi-LLM controller & Yes (trusted, authoritative guideline sources) & Precision/Recall/F1, Consistency (G-Eval), Confusion matrix, Cost analysis & Yes \\
\hline
\end{tabular}
}
\end{table*}


The technical, architectural and evaluation aspects of our agentic AI pipeline differ from existing clinical decision support systems as shown in Table~\ref{tab:comparison_full}. Most systems do not offer formal real-time access to guidelines because they either lack guidelines entirely \cite{r10, r12} or rely on static knowledge bases from single centers \cite{r11, r22}. This limitation prevents them from adapting to new recommendations. Our system retrieves authoritative clinical guidelines in real time and integrates them into the reasoning process. Many solutions focus on narrow domains such as mental health screening with 89\% diagnostic accuracy \cite{r10}, osteoarthritis surgical planning \cite{r11}, or binary heart disease classification with benchmark accuracies over 85\% \cite{r12}. The multi-agent systems ClinicalAgent \cite{r20} and MedGen \cite{r22} use structured reasoning but remain limited to specific clinical trial prediction or static case studies. Our approach explicitly supports both new and follow-up visits across multiple conditions to better match real-world clinical workflows.

The Reasoning Approach and Architecture also show important differences. Simpler ML classifiers and single-task neural networks lack the multi-step reasoning and interpretability required for nuanced decision support. AI-based tools \cite{r14} improve transparency but rely on static single-model pipelines without guideline grounding. The RAG system MEDGPT \cite{r16} enhances evidence use by linking to external documents but remains limited to single-turn and domain-specific question answering. The ClinicalAgent \cite{r20} multi-agent framework achieves better prediction in trial contexts with a PR-AUC of 0.7908 (+0.33 over baseline) but does not incorporate dynamic guidelines. Our pipeline unifies these strengths by combining dynamic retrieval from trusted guidelines with orchestrated, retrieval-augmented multi-LLM reasoning in a controller-based architecture that supports interactive second-opinion generation.

Furthermore, our agentic AI pipeline shows empirical advantages in structured evaluation with strong reasoning-level agreement under G-Eval, matching the highest mean positive-pair score among the evaluated configurations (0.925, tied with Claude Opus 4.6 in Table~\ref{tab:geval_all_cases}). Under the conservative agreement threshold, the agentic configuration achieves precision \(=1.00\), recall \(=0.75\), specificity \(=1.00\), and F1-score \(=0.86\) (Table~\ref{tab:metrics_llms}). It also maintains cost efficiency (average USD~0.0473 per case; Table~\ref{tab:agentic_cost_final}) and practical response times in the tested environment. The results indicate that dynamic retrieval, multi-model orchestration, and guideline-aligned synthesis can provide more reliable second-opinion support than isolated single-model configurations in the evaluated respiratory scenarios.

\subsection[]{Challenges and Limitations}
\label{subsec:limitations}

While the proposed agentic framework demonstrates encouraging performance, its implementation exposed several practical limitations that are important to acknowledge.

A central dependency of the system is the quality and structure of external clinical guidelines. In this study, respiratory disease guidelines from sources such as the World Health Organization and M\'edecins Sans Fronti\`eres were used, and these documents are generally explicit and well organized. This assumption does not hold uniformly across medical domains. In practice, guidelines vary widely in terminology, level of detail, and update frequency. When recommendations are implicit, conditional, or distributed across multiple sections, both retrieval accuracy and downstream reasoning can be affected. There is also a systematic risk that outdated or locally non-applicable guidelines could propagate into recommendations if the underlying corpus is not actively curated and versioned, which places practical constraints on governance and update procedures around guideline ingestion.

The RAG pipeline also introduces sensitivity to how guideline documents are segmented. Although overlapping chunks are used to preserve local context, this approach does not fully address recommendations that span several guideline sections. In such cases, evidence may be retrieved in a fragmented form and must be reconstructed during synthesis. This issue becomes more pronounced for diseases with complex management pathways or recommendations that depend on longitudinal patient history rather than a single clinical snapshot. These behaviors illustrate the brittleness of guideline-based retrieval when key phrases are phrased differently than in the query or scattered across sections, which can lead the system to over-emphasize easily retrievable guidance while under-representing less well-indexed recommendations.

Additional trade-offs arise from the use of an agentic architecture that coordinates multiple LLMs. Combining heterogeneous models improves robustness and reduces reliance on any single system, but it also increases inference time and operational cost. Disagreements between model outputs are not uncommon and require explicit aggregation logic. Even small differences in emphasis or interpretation can influence the final recommendation, which places practical constraints on how aggressively model diversity can be scaled. In the current design, the controller resolves such disagreements by synthesizing a consolidated plan conditioned on the retrieved guideline context and the multiple tool outputs; this reduces variability but can also mask legitimate uncertainty when models disagree for complementary reasons. A more mature deployment would need to expose salient disagreements and the underlying evidence to clinicians instead of collapsing them into a single recommendation, so that inter-model divergence becomes an explicit signal rather than a hidden implementation detail.



Assessing the quality of second medical opinions presents further challenges. The evaluation conducted in this study is limited to a small set of respiratory conditions and relies on an LLM-based evaluator calibrated using a conservative threshold. While this enables scalable comparison across cases, it does not fully reflect inter-clinician variability or the contextual judgment exercised in real clinical practice. Broader evaluation across additional disease domains and assessment by expert clinicians would provide stronger evidence of reliability.

Finally, the system is intentionally positioned as a clinical decision-support tool rather than an autonomous diagnostic system. Its safe use depends on clear communication of uncertainty, explicit grounding in retrieved guideline evidence, and well-defined boundaries regarding responsibility and oversight. These considerations become especially important in higher-risk or complex clinical scenarios, where over-reliance on automated outputs could lead to inappropriate clinical decisions. Taken together, these limitations indicate that rigorous prospective evaluation in clinical settings, including simulation studies and phased roll-outs under supervision, is required before the system could be considered for routine deployment.

\subsection[]{Generalization to Other Disease Domains}
\label{subsec:generalization}

Although the experimental evaluation in this work focuses on respiratory conditions, the underlying system architecture is intended to support extension to other disease domains. As shown in Figure \ref{fig:system_architecture}, the separation between guideline ingestion, retrieval, and agentic reasoning allows generalization to occur primarily through changes to the evidence base rather than through redesign of the orchestration logic.

In practical terms, extending the system to additional diseases requires expanding the guideline corpus used during knowledge base construction. For chronic conditions such as cardiovascular disease or diabetes, this involves incorporating guidelines that are stratified by disease severity, risk categories, and comorbidity profiles. These documents are processed using the same chunking, embedding, and indexing pipeline, allowing retrieval behavior to remain consistent without modification of downstream reasoning components.

More complex disease areas, such as oncology and autoimmune disorders, pose additional challenges. In these domains, clinical recommendations are often conditional on disease stage, prior therapies, and patient-specific biomarkers. Management pathways are therefore highly branched rather than linear. In such cases, the orchestration agent must reason over multiple retrieved guideline fragments that may represent alternative or sequential treatment options. While the current framework supports retrieval and synthesis across multiple chunks, careful aggregation and prompt design are required to avoid loss of clinical context or oversimplified recommendations.

Generalization also requires explicit handling of comorbidity, which is common in many patient populations. Patients with multiple conditions may trigger retrieval from multiple guideline sources, leading to overlapping or potentially conflicting recommendations. Rather than enforcing a single pathway, the proposed system integrates these inputs and can surface areas of agreement or uncertainty. Supporting this form of cross-guideline reasoning is feasible within the existing architecture, but may benefit from additional mechanisms that make trade-offs more explicit to clinicians.

Finally, disease domains with rapidly evolving evidence place increased demands on guideline maintenance. In areas such as oncology, where standards of care change frequently, reliable deployment depends on timely updates to the guideline corpus and clear traceability of retrieved recommendations. These requirements do not change the core architecture, but they highlight the importance of robust data curation and update processes when extending the system beyond relatively stable guideline domains.

\section{CONCLUSION}
\label{sec:conclusion}

In this paper, we presented an agentic AI system for generating evidence-based second medical opinions by combining RAG with controller-based orchestration of multiple LLMs and structured output synthesis. The key contributions of this work include the integration of dynamic clinical guideline retrieval with coordinated multi-model reasoning, the use of an agentic controller to manage tool selection and response synthesis, and the delivery of consistent, guideline-aligned outputs for both new and follow-up clinical cases within a unified framework.

Evaluation within a controlled respiratory-domain setting demonstrates that multi-LLM orchestration can achieve higher reasoning-level agreement with guideline-based reference standards than single-model RAG configurations under a conservative G-Eval threshold. Under the 0.90 agreement threshold, the agentic pipeline achieves precision \(=1.00\), recall \(=0.75\), specificity \(=1.00\), and F1-score \(=0.86\) in the evaluated setting. It also achieves strong semantic alignment with guideline-derived reference answers (mean cosine alignment 0.7629), low per-case token-based costs (average USD~0.0473 across four cases with dynamic tool usage), and mean end-to-end latency of 44.66 seconds in the tested environment. These findings indicate that structured orchestration of heterogeneous LLM tools can improve robustness and reasoning completeness relative to isolated model configurations.

Several issues remain before such systems can be considered for routine clinical use. The current evaluation is limited to simulated respiratory cases and does not establish performance in live clinical environments. System behavior depends on the quality, structure, and currency of ingested guideline sources, and automated alignment metrics do not substitute for clinician-based validation. In addition, inter-model disagreements are currently resolved through controller synthesis rather than being explicitly surfaced to clinicians, which may obscure uncertainty in complex cases.

Future work should therefore focus on four concrete directions required for real clinical value. First, scaling the framework to incorporate broader and more heterogeneous guideline corpora, including cross-specialty and multi-institutional sources, with explicit version control and update governance. Second, extending the input interface beyond structured text to support multi-modal clinical data such as laboratory values, imaging summaries, and longitudinal electronic health record information. Third, conducting clinician-in-the-loop validation studies and prospective pilot deployments to assess usability, trust calibration, and integration with local protocols. Fourth, evaluating system behavior in longitudinal settings that include follow-up visits and longer-term patient trajectories to understand how agentic reasoning performs across evolving clinical contexts.

\section*{CRediT authorship contribution statement}

\textbf{Diana Hawashin}: Formal analysis, Methodology, Investigation, Software, Writing -- original draft. \textbf{Khaled Salah}: Conceptualization, Methodology, Funding acquisition, Project administration, Supervision, Validation, Writing -- review \& editing. \textbf{Raja Jayaraman}: Conceptualization, Validation, Writing -- review \& editing. \textbf{Samer Ellahham}: Conceptualization, Validation, Writing -- review \& editing.


\section*{Acknowledgements}
This research was funded by the Socio-Technical Systems Lab (STSL), Khalifa University of Science and Technology (KU-STSL). 

\section*{Declaration of generative AI and AI-assisted technologies in the writing process}
During the preparation of this work, the authors used ChatGPT in order to improve the readability and language. After using this tool, the authors reviewed and edited the content as needed and take full responsibility for the content of the published article.

\section*{Declaration of Competing Interest}
The authors declare that they have no known competing financial interests or personal relationships that could have appeared to influence the work reported in this paper.

\section*{Data Availability}
The source code and relevant implementation scripts used in this study are publicly available at the following GitHub repository: \url{https://github.com/djh97/smo}.


\section*{Ethical Approval}
This article does not contain any studies with human participants or animals performed by any of the authors. The work is based on simulated datasets and conceptual models. No personally identifiable or protected health information was used.

  

\appendix
\section{Supplementary Material}
\label{app1}

Supplementary material associated with this article, including the released source code, the study case definitions referenced in the evaluation, and the four-case execution trace, can be found in the public repository at \url{https://github.com/djh97/smo}. The representative case execution trace is available at \url{https://github.com/djh97/smo/blob/main/outputs/evaluation/prompt_log_4case.json}, and the study case definitions are available at \url{https://github.com/djh97/smo/blob/main/data/samples/cases.json}.



%% If you have bib database file and want bibtex to generate the
%% bibitems, please use
%%
\bibliographystyle{elsarticle-num} 
\bibliography{refs}

%% End of file `elsarticle-template-num.tex'.
\end{document}
